%==============================================================
% Marmara University - Faculty of Business Administration
% Department of Management Information Systems (B.Sc.)
% Graduation Thesis (May 2026)
% Student:  Abdullah Kılınç (138621049)
% Advisor:  Dr. Öğr. İlkim Ecem Emre
% Compile with pdfLaTeX + biber (Overleaf default works fine).
%==============================================================

\documentclass[12pt,a4paper]{report}

% ---------- Encoding & languages ----------
\usepackage[utf8]{inputenc}
\usepackage[T1]{fontenc}
\usepackage[ngerman,turkish,main=english]{babel}

% ---------- Layout ----------
\usepackage[a4paper,top=2.5cm,bottom=2.5cm,left=2.5cm,right=2.5cm]{geometry}
\usepackage{setspace}
\setstretch{1.35}
\usepackage{indentfirst}
\setlength{\parindent}{1.25cm}
\setlength{\parskip}{0pt}
\usepackage{ragged2e}
\justifying

% ---------- Visuals ----------
\usepackage{graphicx}
\graphicspath{{./}{./figures/}}
\usepackage{amsmath,amssymb,amsfonts}
\usepackage[table]{xcolor}
\definecolor{MarmaraBlue}{HTML}{2F6F9F}
\usepackage{titlesec}
\titleformat{\chapter}[hang]
  {\normalfont\LARGE\bfseries\color{MarmaraBlue}}
  {\thechapter.}{0.6em}{}
\titleformat{\section}
  {\normalfont\Large\bfseries\color{MarmaraBlue}}
  {\thesection}{0.6em}{}
\titleformat{\subsection}
  {\normalfont\large\bfseries\color{MarmaraBlue}}
  {\thesubsection}{0.6em}{}
\titleformat{\subsubsection}
  {\normalfont\normalsize\bfseries\color{MarmaraBlue}}
  {\thesubsubsection}{0.6em}{}
\titlespacing*{\chapter}{0pt}{-20pt}{14pt}

% ---------- Headers/footers ----------
\usepackage{fancyhdr}
\pagestyle{fancy}
\fancyhf{}
\fancyhead[L]{\includegraphics[height=0.85cm]{marmara_isletme_logo.png}}
\fancyfoot[C]{\thepage}
\renewcommand{\headrulewidth}{0pt}
\renewcommand{\footrulewidth}{0pt}
\fancypagestyle{plain}{%
  \fancyhf{}
  \fancyhead[L]{\includegraphics[height=0.85cm]{marmara_isletme_logo.png}}
  \fancyfoot[C]{\thepage}
  \renewcommand{\headrulewidth}{0pt}
}

% ---------- Tables ----------
\usepackage{booktabs}
\usepackage{tabularx}
\usepackage{longtable}
\usepackage{array}
\usepackage{pdflscape}
\usepackage{makecell}

% ---------- Captions ----------
\usepackage[font=small,labelfont=bf,labelsep=period]{caption}

% ---------- Misc ----------
\usepackage{url}
\usepackage{enumitem}
\usepackage{etoolbox}
\usepackage{placeins}
\usepackage{appendix}
\usepackage{listings}

% ---------- Line-breaking & overflow protection ----------
% microtype improves line breaking and reduces overfull boxes;
% emergencystretch + sloppy give TeX more freedom on tight paragraphs;
% the \path{} command (from url) lets long file paths break at "/", "_", ".".
\usepackage{microtype}
\setlength{\emergencystretch}{3em}
\tolerance=1500
\sloppy
% Make \path{} prefer breaks inside long filenames:
\Urlmuskip=0mu plus 1mu
\lstset{
  basicstyle=\ttfamily\footnotesize,
  breaklines=true,
  frame=single,
  columns=fullflexible,
  numbers=none,
  showstringspaces=false
}

% ---------- Hyperlinks ----------
\usepackage[colorlinks=true,linkcolor=MarmaraBlue,citecolor=MarmaraBlue,urlcolor=MarmaraBlue]{hyperref}

% ---------- Bibliography (APA-style author-year, embedded) ----------
% No external .bib file: all references are listed in the thebibliography
% environment at the end of the document. natbib provides \citet and \citep
% in author-year style.
\usepackage[authoryear,round,sort&compress]{natbib}
% Suppress natbib/thebibliography's built-in heading so our \chapter wins:
\renewcommand{\bibsection}{}

%==============================================================
\begin{document}
%==============================================================

\hypersetup{pageanchor=false}

% ============================
% PAGE 1 - GERMAN COVER PAGE
% ============================
\begin{titlepage}
\thispagestyle{empty}
\centering
\vspace*{0.5cm}

{\Large\bfseries MARMARA UNIVERSITÄT}\\[0.4cm]
{\large FAKULTÄT FÜR BETRIEBSWIRTSCHAFT}\\[0.4cm]
{\large WIRTSCHAFTSINFORMATIK (B.SC.)}\\[1.6cm]

{\Large\bfseries\color{MarmaraBlue}
Synthetic Turkey: Eine quellengestützte\\[0.2cm]
generative agentenbasierte Simulation\\[0.2cm]
der türkischen Präsidentschaftswahl 2023\\}
\vspace{0.6cm}
{\small\itshape (Bakkalaureatsarbeit im Studienfach Wirtschaftsinformatik)}

\vfill

\begin{tabular}{rl}
\textbf{vorgelegt von:}        & Abdullah Kılınç \\[0.15cm]
\textbf{Matrikelnummer:}       & 138621049 \\[0.15cm]
\textbf{vorgelegt bei:}        & Dr. Öğr. İlkim Ecem Emre \\[0.15cm]
\textbf{Lehrstuhl:}            & Wirtschaftsinformatik \\
\end{tabular}

\vfill
{\large İSTANBUL · Mai 2026}
\end{titlepage}

\hypersetup{pageanchor=true}

% ============================
% FRONT MATTER (Roman numerals)
% ============================
\pagenumbering{roman}
\setcounter{page}{1}

% ,-- Table of contents ,--
\tableofcontents
\clearpage

% ,-- Summary page ,--
\chapter*{Zusammenfassung / Summary / Özet}
\addcontentsline{toc}{chapter}{Zusammenfassung / Summary / Özet}

\noindent\textbf{TITEL DES PROJEKTS / TITLE OF THE PROJECT / PROJENİN BAŞLIĞI}\\[0.2cm]
Synthetic Turkey: A Source-Grounded Generative Agent-Based Simulation of the 2023 Turkish Presidential Election.\\[0.4cm]

\noindent\textbf{vorgelegt von / submitted by / sunan:} Abdullah Kılınç (Matrikelnummer / student number / öğrenci numarası: 138621049)\\
\textbf{vorgelegt bei / supervised by / danışman:} Dr. Öğr. İlkim Ecem Emre\\[0.4cm]

\noindent\textbf{Zielsetzung – Aim – Amaç:}\\
This thesis investigates whether a source-grounded generative agent-based simulation, in which large language models (LLMs) serve as the reasoning core of synthetic voter agents, can reproduce structurally meaningful patterns of the 2023 Turkish presidential election. The aim is not to predict the result, but to construct a transparent diagnostic instrument that exposes both the promise and the failure modes of LLM-mediated electoral simulation in a non-Western, polarized, two-round presidential context.\\[0.3cm]

\noindent\textbf{Methode – Method – Metot:}\\
A 300-agent synthetic electorate was constructed from twelve research-informed voter archetypes anchored in the 2018 Turkish electoral baseline. Agents were exposed to a longitudinal 2018--2023 event timeline of thirty-seven election-safe ticks. At each tick, an LLM (\texttt{gpt-4o-mini}, temperature 0.45) updated voter beliefs, affect, trust, turnout probability, candidate intention, and a short first-person reflection through a strict JSON output schema. Results were evaluated at five levels: aggregate vote share, archetype-level voting, demographic breakdowns, longitudinal trust/affect trajectories, and tie-handling sensitivity.\\[0.3cm]

\noindent\textbf{Material / Datensatz – Material / Data set – Materyal / Veri seti:}\\
The dataset is fully synthetic and comprises (i) 300 voter-agent persona files (JSON), (ii) twelve archetype configuration files (JSON/YAML) with worldview, media-diet, behavioural and affective parameters, (iii) the 2018 baseline sampling profile (YAML), (iv) political-actor and credibility-matrix configuration files (YAML), (v) a 37-tick event timeline (JSON) covering the 2018--2023 Turkish political period, (vi) the exported long-format agent-trajectory table (\texttt{agent\_trajectories.csv}, 11{,}100 rows, 35 columns), (vii) per-agent belief and episodic memory archives (300 belief JSON files plus 300 episodic JSON files), and (viii) the canonical run log (\path{run_2023_20260516_142653_179132.jsonl}, 11{,}138 records).\\[0.3cm]

\noindent\textbf{Verwendetes Tool – Used tool – Kullanılan araç:}\\
Python 3.11; the OpenAI API (\texttt{gpt-4o-mini}, temperature 0.45, 900 output-token budget) for voter-decision and political-broadcast generation; PyYAML and JSON for configuration loading; pandas and NumPy for analysis; matplotlib and seaborn for visualization; Python's \texttt{concurrent.futures.ThreadPoolExecutor} for parallel voter calls.\\[0.3cm]

\noindent\textbf{Stichwörter – Keywords – Anahtar kelimeler:}\\
generative agent-based modelling; large language models; synthetic voters; electoral simulation; Turkish politics.

\clearpage

% ,-- List of Tables ,--
\listoftables
\addcontentsline{toc}{chapter}{List of Tables}
\clearpage

% ,-- List of Figures ,--
\listoffigures
\addcontentsline{toc}{chapter}{List of Figures}
\clearpage

% ,-- List of Abbreviations ,--
\chapter*{List of Abbreviations}
\addcontentsline{toc}{chapter}{List of Abbreviations}

\begin{tabularx}{\linewidth}{@{}lX@{}}
\toprule
\textbf{Abbreviation} & \textbf{Expansion} \\
\midrule
ABM & Agent-Based Model / Modelling \\
AKP & Adalet ve Kalkınma Partisi (Justice and Development Party) \\
ANES & American National Election Studies \\
ATA & ATA İttifakı (ATA Alliance) \\
CHP & Cumhuriyet Halk Partisi (Republican People's Party) \\
CI  & Confidence Interval \\
CSS & Computational Social Science \\
DEM & Demokratik Halklar Partisi (Peoples' Democratic Party, post-2023 successor) \\
GABM & Generative Agent-Based Modelling \\
HDP & Halkların Demokratik Partisi (Peoples' Democratic Party) \\
JCSS & Journal of Computational Social Science \\
JSON & JavaScript Object Notation \\
LLM & Large Language Model \\
MHP & Milliyetçi Hareket Partisi (Nationalist Movement Party) \\
MIS & Management Information Systems \\
RAG & Retrieval-Augmented Generation \\
T030A & First-Round Vote-Decision Tick \\
T035A & Runoff Vote-Decision Tick \\
YRP & Yeniden Refah Partisi (New Welfare Party) \\
YAML & YAML Ain't Markup Language \\
YSK & Yüksek Seçim Kurulu (Supreme Election Council of Türkiye) \\
YSP & Yeşil Sol Parti (Green Left Party) \\
\bottomrule
\end{tabularx}
\clearpage

% ,-- List of Appendices ,--
\chapter*{List of Appendices}
\addcontentsline{toc}{chapter}{List of Appendices}

\begin{tabularx}{\linewidth}{@{}lX@{}}
\toprule
\textbf{Appendix} & \textbf{Title} \\
\midrule
A & Voter Archetype Profiles \\
B & Prompt Template and JSON Output Schema \\
C & Simulation Event Timeline (37 Ticks) \\
D & Tie-Handling Sensitivity Detail \\
\bottomrule
\end{tabularx}
\clearpage

% ============================
% MAIN MATTER (Arabic numerals)
% ============================
\pagenumbering{arabic}
\setcounter{page}{1}

%==============================================================
\chapter{Einführung -- Introduction}
\label{ch:introduction}
%==============================================================

\section{Zweck und Umfang des Projekts -- Purpose and Scope of the Project}
\label{sec:scope}

This thesis presents a computational study of voter behaviour in the 2023 Turkish presidential election through a method known as \emph{generative agent-based modelling} (GABM). Generative agent-based modelling is a recent extension of traditional agent-based modelling in which the reasoning component of each simulated agent is provided by a large language model (LLM) rather than by a hand-coded rule set. Each agent is given a persona, a memory, an environment, and a stream of events, and the LLM is asked to respond in character. The system as a whole therefore consists of two layers. The deterministic outer layer organises the population, defines the events, controls the information flow, validates outputs, and aggregates results. The probabilistic inner layer, the LLM, performs the cognitive step of interpreting each event from the perspective of an individual voter.

The project, named \emph{Synthetic Turkey}, applies this paradigm to the May 2023 Turkish presidential election. A synthetic electorate of three hundred voter agents is constructed from twelve research-informed archetypes that are calibrated to the 2018 Turkish parliamentary and presidential baseline. The agents experience a longitudinal stream of thirty-seven simulation ticks reflecting the major political, economic, and social events that took place between June 2018 and May 2023. At each tick the LLM produces, for each agent, a structured update of trust, government approval, emotional state, party preference, candidate-level probabilities, turnout intention, and a short first-person reflection. The simulation closes with election-safe decision ticks at which the simulated first round and runoff are recorded before the actual electoral outcome is revealed inside the simulation.

The aim of the project is methodological rather than predictive. The Turkish 2023 election is a particularly demanding case for any electoral simulation. It took place in a polarized two-round presidential system, against a backdrop of currency crisis and rapid inflation, sustained democratic backsliding, opposition alliance formation, Kurdish party repression and reorganization, the February 2023 Kahramanmaraş earthquakes, nationalist fragmentation that produced the third candidate Sinan Oğan, and pre-election polling which mostly forecast an opposition victory that did not materialize. A computational model that can reproduce a result of this kind by accident is not necessarily a model that captures the underlying mechanisms. A model that fails to reproduce it, by contrast, may still offer interpretable information about which assumptions broke and why.

For these reasons, the present thesis frames its three research questions narrowly. \textbf{Research Question 1} asks: \emph{to what extent can a source-grounded LLM-based agent-based simulation reproduce the aggregate outcome of the 2023 Turkish presidential election under different probability-to-vote conversion rules?} \textbf{Research Question 2} asks: \emph{at which levels of analysis (aggregate, archetype, demographic, and longitudinal) does the simulation succeed or fail to reproduce empirically observed electoral patterns?} \textbf{Research Question 3} asks: \emph{what diagnostic lessons does the simulation offer for the design of LLM-mediated electoral simulation methodology, particularly regarding aggregation rules, calibration, archetype coverage, and the conversion of probabilistic intention into discrete electoral choice?}

The study makes three corresponding claims. First, it offers a transparent and reproducible pipeline for building a source-grounded synthetic electorate in a non-Western political context. Second, it evaluates that pipeline at five distinct levels of analysis, so that aggregate-level accuracy is never used as the sole measure of success. Third, it identifies and discusses six concrete failure modes that are likely to affect any future LLM-agent electoral simulation: sensitivity to aggregation rules, sensitivity to tie handling, omission of protest-vote pathways, over-homogenization of certain archetypes, latent LLM knowledge of the election outcome, and weak numerical responsiveness in structured state variables \citep{kang2026dynamok,qu2024performance,rettenberger2025politicalbias}.

What the thesis explicitly does \emph{not} claim is also important. It does not claim that LLM agents are an empirical substitute for polling. It does not claim that the simulation captures every causal pathway in Turkish voting behaviour. It does not claim that any individual agent corresponds to any real voter. The single 300-agent run reported in this thesis is a stochastic realisation of a model that, on the basis of one run alone, can be audited but cannot be statistically validated in the strict sense. The contribution of the thesis is therefore framed as a diagnostic methodology study: a controlled demonstration of what an LLM electorate built on best-practice principles can and cannot reproduce, and what those results imply for future research \citep{kreutner2026ep,qu2024performance}.

%==============================================================
\chapter{Theoretischer Teil -- Theoretical Part}
\label{ch:theory}
%==============================================================

\section{Agent-Based Modelling}
\label{sec:abm}

Agent-based modelling is a class of computational methods in which a population of autonomous entities, the agents, interact with each other and with their environment according to specified rules. Each agent has a state (for example, a political opinion, an income level, or a willingness to participate) and a behaviour function that decides how it responds to stimuli and how its state evolves over time. The macro-level pattern of the simulation is not programmed directly. Instead, it emerges from the interactions of many heterogeneous agents acting locally on local information \citep{epstein2006generative,macal2010tutorial}.

Two properties make agent-based modelling attractive for the study of political behaviour. First, the method handles heterogeneity gracefully: each agent may differ from every other agent in identity, history, geography, and disposition, which is closer to the reality of electoral populations than mean-field approaches. Second, the method allows the modeller to separate \emph{micro-foundations} (what individual agents do) from \emph{macro-outcomes} (what the system as a whole produces), which makes it well suited to questions about how individual reasoning generates collective political phenomena.

Traditional agent-based models, however, encode agent reasoning in rule sets or utility functions that are written by the modeller in advance. This limits the expressive range of agent behaviour and forces the modeller to commit to a specific theory of cognition before the simulation begins. Recent work has therefore turned to large language models as a more flexible behavioural layer for individual agents \citep{gao2024llmabm,lu2024gabs}.

\section{Large Language Models}
\label{sec:llms}

Large language models are statistical models trained on very large corpora of natural-language text. They generate output by predicting, one token at a time, what continuation is most plausible given the preceding context. Through that training they acquire a broad latent representation of natural-language usage, including, often by accident, considerable knowledge of historical events, public figures, and cultural reference points. They are not, however, calibrated probability machines: their outputs reflect what is plausible given training rather than what is statistically grounded in a population \citep{argyle2023out,aher2023simulating,qu2024performance}.

A central question in computational social science is whether LLMs can play the role of \emph{simulated respondents}: agents that approximate the distribution of human survey responses for a given population. Early work suggested that conditioning an LLM on a demographic and ideological profile can recover some of the patterns observed in survey data \citep{argyle2023out}. More recent work has argued that profile-based conditioning is fragile: prediction accuracy improves substantially when models are also conditioned on relevant past opinions, suggesting that demographics alone are not enough to determine an LLM's response \citep{hwang2023aligning}. There is also a growing literature documenting systematic ideological biases in LLM outputs and cautioning against treating LLM-generated probabilities as calibrated empirical probabilities \citep{rettenberger2025politicalbias,qu2024performance,ziems2024css}.

These considerations matter for the present thesis. The synthetic voters are conditioned on persona, archetype, prior opinion, memory, and event context (not on a static demographic vector alone), precisely because the literature warns that the LLM's behaviour is unstable when conditioning is too thin.

\section{Generative Agent-Based Modelling}
\label{sec:gabm}

Generative agent-based modelling combines the population-level structure of traditional ABM with the open-ended reasoning capacity of LLMs. In the architecture introduced by \citet{park2023generative}, each agent maintains a stream of memories, retrieves relevant past experiences when responding to new stimuli, periodically generates higher-level reflections from clusters of memories, and uses both memories and reflections to plan future behaviour. The result is a simulation in which the agent's text-level behaviour appears continuous and believable rather than rule-driven.

Subsequent work has positioned GABM as an emerging paradigm for the study of complex social systems, while at the same time emphasizing four unresolved methodological problems: prompt sensitivity (small changes in prompt wording can change agent behaviour), hallucination (agents may produce content that is internally inconsistent or factually wrong), model-specific bias (different LLMs produce systematically different distributions of responses), and standardization (the field has not yet agreed on common protocols for reporting GABM studies) \citep{gao2024llmabm,lu2024gabs}.

The Synthetic Turkey project sits inside this paradigm but with two deliberate restrictions. First, the simulation is bounded by a historical timeline rather than open-ended: agents respond to specific dated events, not to free-form social interaction. Second, the deterministic outer layer (event order, population sampling, broadcast filtering, JSON validation, and aggregation rules) is held strictly outside the LLM. This separation is intended to make the simulation auditable: the share of the result that can be attributed to the LLM is, in principle, isolable from the share that is determined by the configuration.

\section{Electoral Behaviour and Key Concepts}
\label{sec:electoral_concepts}

The thesis uses a small number of concepts from political behaviour research that need to be defined for the non-specialist reader. \emph{Retrospective voting} is the idea that voters evaluate incumbents on the basis of past performance, especially in economic and disaster-response domains. \emph{Affective polarization} is the tendency for partisans to feel emotional attachment to their own side and emotional hostility toward the opposing side, often more strongly than they feel ideological commitment to specific policies \citep{aytac2025polarised}. \emph{Strategic voting} is the practice of casting a ballot for a candidate other than one's most preferred candidate in order to maximise the probability of a desired electoral outcome. \emph{Cross-pressure} describes the condition of a voter whose various identities or interests pull in opposing political directions, for example a culturally conservative voter under acute economic stress.

These concepts matter for the design of the simulated archetypes. The Earthquake Zone Loyalist archetype, for example, embodies a non-trivial finding from the retrospective-voting literature: that the February 2023 earthquakes did not produce automatic punishment of the incumbent in affected provinces, because voters appear to have weighted aid delivery and political identity alongside disaster-response performance \citep{sevinc2024afet}. The Conservative Economically Disillusioned archetype, by contrast, embodies the cross-pressure condition: voters with deep economic grievances but enduring cultural and religious attachments to the incumbent coalition \citep{carkoglu2024unlikely,yildirim2025sege}. Affective polarisation appears in the Devout Anatolian Loyalist and Kurdish Political Voter archetypes, on opposite sides of the affective cleavage \citep{aytac2025polarised,baruh2024difficult}.

\section{The 2023 Turkish Presidential Election Context}
\label{sec:context}

The 2023 Turkish presidential election was held in two rounds, on 14 May and 28 May 2023. In the first round Recep Tayyip Erdoğan obtained 49.52 per cent, his main challenger Kemal Kılıçdaroğlu obtained 44.88 per cent, the nationalist independent Sinan Oğan obtained 5.17 per cent, and the late-stage candidate Muharrem İnce obtained 0.43 per cent. Because no candidate exceeded the fifty-per-cent threshold, a runoff was held between Erdoğan and Kılıçdaroğlu on 28 May 2023, which Erdoğan won with 52.18 per cent against 47.82 per cent \citep{ysk2023}.

The election took place in an unusually demanding political environment. The Turkish economy had experienced a major currency crisis in 2018, sustained high inflation through 2021 and 2022, and a sharp decline in real wages and pensions. The February 2023 earthquakes affected eleven provinces and ten million people, and triggered an extended public debate about state capacity and the government response \citep{sevinc2024afet}. The opposition formed a Nation Alliance under the leadership of Kılıçdaroğlu and faced repeated dilemmas about candidate selection, including substantial centre-right and İYİ Party reluctance to support Kılıçdaroğlu personally \citep{baruh2024difficult}. The Kurdish political movement, organized at that time under the HDP/YSP umbrella, faced an active closure case and chose strategic coordination with Kılıçdaroğlu rather than a separate presidential candidacy \citep{ilbiz2023securitization,yegen2022kurdish}. Anti-immigration sentiment, partly mobilised by the newly visible Victory Party, contributed to a third-candidate pathway around Sinan Oğan \citep{gerim2025victory}. Polarisation across these cleavages was strong and shaped both turnout and inter-bloc transfers \citep{aytac2025polarised,carkoglu2024unlikely,cinkara2025erdogan}.

For the purposes of an LLM-agent simulation, what makes the 2023 election demanding is not the existence of any single cleavage but the way these cleavages combined. A model that simulates only the religious/secular dimension, or only economic discontent, will not capture the result. A model that aggregates probabilistically without considering how voters actually convert their preferences into a single ballot risks producing an outcome that is internally coherent but empirically wrong \citep{kang2026dynamok,kreutner2026ep}.

%==============================================================
\chapter{Literaturrecherche -- Literature Review}
\label{ch:litreview}
%==============================================================

\section{Search Strategy, Databases, and Keywords}
\label{sec:search_strategy}

The literature reviewed for this thesis was assembled through structured searches in three academic databases: \emph{Google Scholar}, \emph{Scopus}, and \emph{DergiPark}. Google Scholar was used as the primary discovery tool because it indexes the broadest range of English-language preprints, conference proceedings, and peer-reviewed journal articles relevant to large language model research and agent-based modelling. Scopus was used to verify the indexing status of the major peer-reviewed sources and to identify English-language empirical studies of the 2023 Turkish elections in indexed outlets such as the \emph{Journal of Common Market Studies}, \emph{Democratization}, and \emph{South European Society and Politics}. DergiPark, the central index of Turkish academic journals, was used to locate Turkish-language empirical studies of the 2023 election, earthquake-related voting, party-system change, and youth political behaviour, since these works are not consistently indexed in Scopus or Google Scholar.

The searches were performed using nine keyword groups, in both English and Turkish where applicable. The first group covered the methodological foundation (\enquote{agent-based modelling}, \enquote{agent-based simulation}). The second group covered its recent extension (\enquote{generative agent-based modelling}, \enquote{generative agents}, \enquote{LLM agents}). The third group covered large language models as social-science instruments (\enquote{large language models}, \enquote{LLM simulation}, \enquote{silicon samples}). The fourth group targeted electoral applications (\enquote{synthetic voters}, \enquote{electoral simulation}, \enquote{voter simulation}). The fifth group addressed methodology (\enquote{computational social science}). The sixth group targeted the Turkish case (\enquote{Turkish elections}, \enquote{2023 Turkish presidential election}, \enquote{T\"urkiye 2023 cumhurba\c{s}kanl\i\u{g}\i\ se\c{c}imi}, \enquote{se\c{c}men davran\i\c{s}\i}). The seventh group covered case-specific terms (\enquote{Kahramanmara\c{s} depremi se\c{c}men davran\i\c{s}\i}, \enquote{Sinan O\u{g}an}, \enquote{ATA \.Ittifak\i}, \enquote{\.IY\.I Parti}, \enquote{HDP / YSP / DEM}, \enquote{Yeniden Refah Partisi}). The eighth group targeted partisan psychology (\enquote{affective polarisation}, \enquote{duygusal kutupla\c{s}ma}). The ninth group covered methodological risk (\enquote{political bias in LLMs}, \enquote{LLM political opinion}).

The cross-database searches returned more than four hundred candidate references after duplicate removal. The candidate pool was filtered to retain only peer-reviewed publications, accepted conference proceedings, recognised preprint archives, academic book chapters, and authoritative data sources; blog posts and untraceable working papers were excluded. After this filter approximately one hundred and twenty references remained relevant. On the basis of substantive relevance to either (a) LLM-agent simulation methodology or (b) the Turkish 2023 case, \textbf{thirty-two studies were retained for detailed examination}. Each is described in the prose paragraphs that follow, organised into six thematic clusters; the database in which the study was discovered and the primary keyword that surfaced it are recorded at the start of each entry, together with the study code (ST01 -- ST32) used for cross-reference inside this thesis.

\section{Cluster 1: Agent-Based Modelling and Generative Agent-Based Modelling}
\label{sec:litcluster1}

The conceptual foundation of the thesis is supplied by five studies that span six decades of agent-based modelling and its recent extension into LLM-driven generative agents. The earliest is \citet{epstein2006generative} (\textbf{ST01}; discovered in \emph{Google Scholar} using the keyword \enquote{agent-based modelling}). Epstein's book \emph{Generative Social Science} positions agent-based modelling at the centre of what he calls \emph{generative social science}, the position that a social-science model has explanatory value only when the macro-pattern of interest can be \emph{generated} from heterogeneous micro-level rules, not merely described. This conceptual stance frames the present thesis as a generative social-science project rather than as a curve-fitting exercise.

A complementary didactic foundation is provided by \citet{macal2010tutorial} (\textbf{ST02}; \emph{Google Scholar}; keyword \enquote{agent-based modelling}). The authors offer a step-by-step tutorial that defines an agent, distinguishes ABM from other simulation paradigms, and reviews the practical stages of model design, calibration, and validation. The paper is the canonical introductory reference in the ABM literature and is used in this thesis to define basic vocabulary for the non-specialist committee reader.

The leap from rule-based ABM to LLM-driven ABM is anchored in \citet{park2023generative} (\textbf{ST03}; \emph{Google Scholar}; keyword \enquote{generative agents}). Park and colleagues introduced an architecture in which each LLM-driven agent stores memories, retrieves them by relevance, generates higher-level reflections, and plans future behaviour. Although the original demonstration was a small simulated village rather than an electoral simulation, the architectural pattern of memory--reflection--planning is directly reused in Synthetic Turkey, where each voter agent maintains a persistent belief state, an episodic memory of past ticks, and a per-tick first-person reflection.

The methodological landscape that has emerged around generative agents is mapped by two recent surveys. \citet{gao2024llmabm} (\textbf{ST04}; \emph{Scopus}; keyword \enquote{large language models empowered agent-based modelling}) offer a structured taxonomy of the LLM-ABM literature and explicitly identify four unresolved problems that any GABM study must address: prompt sensitivity, hallucination, model-specific ideological bias, and the absence of standardised reporting protocols. Each of these concerns is engaged explicitly in the methods and limitations chapters of the present thesis. \citet{lu2024gabs} (\textbf{ST05}; \emph{Scopus}; keyword \enquote{generative agent-based models}) provide a conceptual review for complex-systems research and argue that the deterministic environment around the LLM must be carefully separated from the LLM itself. Synthetic Turkey follows this guidance by keeping event order, population sampling, broadcast filtering, JSON validation, and aggregation strictly outside the LLM, so that the LLM is responsible only for the voter-reasoning step.

\section{Cluster 2: LLMs as Synthetic Respondents and Memory-Conditioned Agents}
\label{sec:litcluster2}

A second cluster of studies treats LLMs not as components of a sandbox simulation but as stand-ins for individual respondents in survey-like settings. The defining reference is \citet{argyle2023out} (\textbf{ST06}; \emph{Google Scholar}; keyword \enquote{silicon samples}), who propose the \emph{silicon samples} framing: an LLM conditioned on a demographic and ideological profile is asked the same items that a human respondent would answer in a public-opinion survey. The authors report encouraging agreement between the simulated and observed American National Election Studies distributions. The framing is foundational for any LLM-based voter simulation; the present thesis adopts it but layers memory, event exposure, and longitudinal state on top, in line with later critiques.

\citet{aher2023simulating} (\textbf{ST07}; \emph{Google Scholar}; keyword \enquote{simulate human subject studies}) extend the silicon-sample idea by testing whether LLMs can replicate the findings of classic human-subject experiments in psychology. Their results are mixed: some experiments replicate, others do not. The paper provides early empirical evidence that LLM behaviour is task-dependent and cannot be assumed to mirror human behaviour by default.

The most direct design implication for the present project comes from \citet{hwang2023aligning} (\textbf{ST08}; \emph{Google Scholar}; keyword \enquote{aligning language models to user opinions}). The authors show that conditioning an LLM on a demographic and ideological profile alone is insufficient to predict an individual user's opinions; adding relevant past opinions to the prompt produces systematic accuracy gains. This finding has direct architectural consequences. It implies that LLM-based voter agents conditioned on only a demographic vector are likely to be miscalibrated, while agents that also receive prior political memory and an evolving state are likely to behave more plausibly. Synthetic Turkey follows this guidance closely: each voter agent carries a 2018 baseline memory, a persona biography, accumulated episodic memories, and an evolving belief and affect state rather than a single static archetype label.

A broader cautionary voice is provided by \citet{ziems2024css} (\textbf{ST09}; \emph{Scopus}; keyword \enquote{computational social science large language models}). Their evaluation of thirteen LLMs across twenty-five computational social science benchmarks concludes that LLMs can augment social-science workflows, especially in annotation and explanation tasks, but that their performance remains task-dependent and does not remove the need for traditional validation. The conclusion is consistent with the framing of the present thesis as a diagnostic instrument rather than as a validated predictive tool.

\section{Cluster 3: LLM-Based Political and Electoral Simulation}
\label{sec:litcluster3}

A third cluster contains studies that apply LLM agents directly to political or electoral settings. \citet{zhou2025flockvote} (\textbf{ST10}; \emph{Google Scholar}; keyword \enquote{electoral simulation large language model}) present FlockVote, a system that simulates the 2024 U.S. presidential election using LLM agents with demographic profiles and dynamic contextual information. \citet{zhang2024electionsim} (\textbf{ST11}; \emph{Google Scholar}; keyword \enquote{ElectionSim}) take a more data-intensive approach, constructing a voter pool of approximately one million LLM-driven agents from social-media data and reporting strong state-level accuracy in U.S. presidential simulation. Together, these studies show that LLM-agent electoral simulation is technically feasible at scale and that aggregate-level patterns in the U.S. case can be matched. They are, however, both anchored in a two-party, single-round U.S. setting that differs structurally from the 2023 Turkish election.

\citet{kang2026dynamok} (\textbf{ST12}; \emph{Google Scholar}; keyword \enquote{Korean LLM political bias}) reports a census-grounded Korean-language electoral simulation that diagnoses three systematic failure modes: progressive bias in LLM outputs, regional-polarisation collapse, and the collapse of third-party salience into major-bloc reasoning. The third pattern, third-party collapse, directly parallels the Sinan O\u{g}an miss identified in the present thesis and is the strongest non-Western precedent for the diagnostic-failure framing developed here.

\citet{kreutner2026ep} (\textbf{ST13}; \emph{Google Scholar}; keyword \enquote{persona-driven voting LLM}) test persona-driven LLM simulation against European Parliament voting behaviour. The authors find that persona conditioning can approximate observed roll-call behaviour but report substantial sensitivity to prompt wording, counterfactual setup, and generation method. The result supports the present thesis's design choice of evaluating Synthetic Turkey not only by final winner accuracy but also by coalition structure, archetype-level behaviour, and aggregation-rule sensitivity.

\section{Cluster 4: Validity Risks, Bias, and Calibration}
\label{sec:litcluster4}

A fourth cluster addresses the validity risks of LLM-based simulation directly. \citet{rettenberger2025politicalbias} (\textbf{ST14}; \emph{Scopus}; keyword \enquote{political bias in large language models}) measure systematic ideological alignment patterns across multiple LLMs and conclude that LLMs are not ideologically neutral instruments. The paper is published in the \emph{Journal of Computational Social Science} (the same journal targeted by the companion manuscript of this thesis), and provides a direct precedent for treating LLM ideological priors as a first-order validity concern. \citet{qu2024performance} (\textbf{ST15}; \emph{Scopus}; keyword \enquote{public opinion simulation LLM}) provide cross-national evidence that LLM performance in public-opinion simulation varies substantially by country, demographic group, and topic, with particular weakness in non-Western and non-English contexts. Together with \citet{kang2026dynamok}, this work argues that apparent winner accuracy in an LLM-agent electoral simulation can mask substantial errors in vote share and in third-party representation. The present thesis responds to these warnings by reporting outcomes under multiple aggregation rules, by examining coalition-level behaviour explicitly, and by presenting tie-handling sensitivity as a substantive finding rather than a reporting detail.

\section{Cluster 5: Turkish Electoral Context and the 2023 Election}
\label{sec:litcluster5}

The largest cluster contains the empirical literature on the 2023 Turkish election. \citet{carkoglu2024unlikely} (\textbf{ST16}; \emph{Scopus}; keyword \enquote{Turkey 2023 elections}) argues that Erdo\u{g}an's 2023 survival cannot be explained by economic evaluations alone. Conservative moral values, security, foreign policy, religiosity, right-wing identity, and regime satisfaction are identified as decisive contributors; Kurdish and Alevi voters were negatively predisposed toward Erdo\u{g}an. This paper is the single most cited empirical reference in the present thesis. It supports the construction of the Devout Anatolian Loyalist, Conservative Economically Disillusioned, Alevi-CHP Loyalist, and Kurdish Political Voter archetypes used in the simulation.

\citet{aytac2025polarised} (\textbf{ST17}; \emph{Scopus}; keyword \enquote{affective polarisation Turkey 2023}) report a two-wave panel-survey analysis showing that affective polarisation was particularly strong on the AKP and HDP poles and that polarisation reinforced AKP voter loyalty across the campaign. The finding underpins the modelling choice to treat partisan affect and loyalty as central state variables in the voter prompt rather than reducing voters to purely economic actors.

\citet{baruh2024difficult} (\textbf{ST18}; \emph{Scopus}; keyword \enquote{Nation Alliance 2023 candidate choice}) examine candidate-choice tradeoffs inside the opposition Nation Alliance. The authors show that K\i l\i\c{c}daro\u{g}lu held a comparative advantage among HDP-affinity voters but a comparative disadvantage among \.IY\.I-aligned centre-right voters. This result supports the simulation's representation of the Kurdish Political Voter archetype as a strategic anti-incumbent voter rather than as a conventional CHP voter, and it underpins the Moderate Nationalist \.IY\.I archetype's mixed disposition.

\citet{ilbiz2023securitization} (\textbf{ST19}; \emph{Scopus}; keyword \enquote{competitive authoritarianism opposition alliance}) provide the institutional context. Using a process-tracing approach, the authors show how securitisation and fear politics shaped opposition coordination in competitive-authoritarian regimes, including Turkey. The paper informs the simulation's treatment of strategic Kurdish coordination behind the Nation Alliance under securitisation pressure.

\citet{sevinc2024afet} (\textbf{ST20}; \emph{DergiPark}; keyword \enquote{Kahramanmara\c{s} depremi se\c{c}men davran\i\c{s}\i}) report a field-survey analysis of voter behaviour in the earthquake-affected provinces of February 2023. They argue that affected voters engaged in \emph{mediated retrospective voting}, weighing aid delivery, services, and political identity alongside disaster-response performance, rather than mechanically punishing the incumbent. This is the empirical anchor for the Earthquake Zone Loyalist archetype.

\citet{zeynel2024duygular} (\textbf{ST21}; \emph{DergiPark}; keyword \enquote{politik duygular se\c{c}im kampanyas\i}) provide a qualitative content analysis of the most-watched campaign advertisements of the six highest-vote-share parties in the 2023 election. The paper shows that emotional framing differed systematically by party ideology and by ruling/opposition position. This supports the simulation's treatment of political broadcasts as affective and partisan inputs rather than as neutral information packets.

\citet{gerim2025victory} (\textbf{ST22}; \emph{DergiPark}; keyword \enquote{Zafer Partisi g\"o\c{c}}) analyses the influence of anti-immigration attitudes on voter behaviour and the case of the Victory Party. The paper documents the real but bounded protest-vote audience that the simulation's missing nationalist-protest archetype would need to capture in order to reproduce the Sinan O\u{g}an vote share; this is one of the most direct empirical grounds for the diagnostic finding of the present thesis.

\citet{karaca2024socialmedia} (\textbf{ST23}; \emph{DergiPark}; keyword \enquote{2023 cumhurba\c{s}kanl\i\u{g}\i\ sosyal medya}) reports a content analysis of the Instagram and X (formerly Twitter) use of the 2023 presidential candidates. The paper documents the role of social-media-mediated political communication for younger voters and supports the inclusion of social-media exposure as a separate media-diet channel in the voter configuration.

\citet{yildirim2025sege} (\textbf{ST24}; \emph{DergiPark}; keyword \enquote{sosyo-ekonomik geli\c{s}mi\c{s}lik se\c{c}men tercihi}) provides an aggregate-level regression analysis of socio-economic development and voter preferences across the eighty-one Turkish provinces for the 28 May 2023 runoff. The paper supports the regional differentiation embedded in the synthetic electorate.

\citet{cinkara2025erdogan} (\textbf{ST25}; \emph{DergiPark}; keyword \enquote{Erdo\u{g}an se\c{c}im ba\c{s}ar\i s\i\ yeni milliyet\c{c}ilik}) supply a qualitative analytical essay on Erdo\u{g}an's electoral success that emphasises geopolitics, social class, and \emph{new nationalism} in coalition consolidation. The paper informs the design of the Devout Anatolian Loyalist and Conservative Economically Disillusioned archetypes.

\citet{bozdag2022skeptical} (\textbf{ST26}; \emph{Scopus}; keyword \enquote{misinformation Turkey polarisation}) provide mixed-methods evidence on news consumption and misinformation in Turkey, documenting a state of \emph{skeptical inertia} among polarised media consumers. The paper supports the design of the credibility matrix and the asymmetric exposure of pro-government versus opposition broadcasts in the simulation.

\citet{najafi2024dataset} (\textbf{ST27}; \emph{Scopus}; keyword \enquote{2023 Turkish election dataset}) released the first public social-media dataset for the 2023 Turkish general election, comprising a large tweet collection with account, timestamp, and content metadata. The dataset is cited here as a documentation resource for the broader Turkish 2023 social-media environment; the present thesis does not consume it directly but acknowledges it as a contemporaneous data infrastructure.

\section{Cluster 6: Background Studies on Turkish Cleavages and Generational Politics}
\label{sec:litcluster6}

A sixth cluster contains the broader background literature on Turkish party politics and generational behaviour. \citet{ozbudun2013cleavages} (\textbf{ST28}; \emph{Google Scholar}; keyword \enquote{Turkish party politics social cleavages}) supplies the classical reference work on Turkish party politics and social cleavages. The book frames the long-run secular/religious, ethnic, and centre/periphery divisions that the simulation's archetypes are designed to represent.

\citet{kuru2017religious} (\textbf{ST29}; \emph{Google Scholar}; keyword \enquote{Turkey constitution religion}) examines constitution-writing and religious divisions in Turkey. The book chapter supplies the comparative-constitutional context for the secular--religious cleavage that underpins the Alevi-CHP Loyalist and Devout Anatolian Loyalist archetypes.

\citet{yegen2022kurdish} (\textbf{ST30}; \emph{Google Scholar}; keyword \enquote{Erdo\u{g}an Kurdish question}) provide a policy commentary on Erdo\u{g}an and the Turkish opposition's revisiting of the Kurdish question on the eve of the 2023 election. The paper informs the simulation's treatment of strategic Kurdish coordination behind the Nation Alliance candidate and the operation of the Kurdish movement state machine in the political-actor layer.

\citet{demir2024iyi} (\textbf{ST31}; \emph{DergiPark}; keyword \enquote{\.IY\.I Parti}) supplies a theoretical evaluation of the \.IY\.I Party within the framework of political-party theory. The paper is used to ground the Moderate Nationalist \.IY\.I archetype in the centre-right nationalist tradition rather than treating it as a residual category.

\citet{hafizoglu2021z} (\textbf{ST32}; \emph{DergiPark}; keyword \enquote{Z ku\c{s}a\u{g}\i\ se\c{c}men}) examines the voting behaviour of Generation~Z in Turkey, documenting lower traditional party loyalty, higher social-media exposure, and a stronger opposition lean among younger voters. The paper is the principal source for the Young Urban Protest Voter archetype.

\section{The Gap That Synthetic Turkey Fills}
\label{sec:gap}

Read together, the thirty-two reviewed studies point to a clear methodological gap. The LLM-agent simulation literature (Clusters~1--4) has produced impressive demonstrations of synthetic respondents and electoral simulation in the United States and has begun to test the architecture in Korea (\textbf{ST12}) and in the European Parliament (\textbf{ST13}). The Turkish electoral literature (Clusters~5--6) has produced rich empirical and theoretical accounts of the 2023 election and the cleavages that structure Turkish politics. Yet the two literatures have not been combined. No prior study has, to the best of this author's knowledge, built a source-grounded LLM-agent electoral simulation for a non-Western two-round presidential election with the level of structural complexity exhibited by T\"urkiye in 2023. Synthetic Turkey is positioned at the intersection of these two literatures: it uses Clusters~5 and~6 as its empirical scaffolding and applies the methodological apparatus of Clusters~1 and~3 to that scaffolding, while remaining explicitly attentive to the validity warnings of Cluster~4. The specific contribution of the thesis is to test whether an LLM-agent simulation, when grounded in this rich empirical literature, can reproduce the structural patterns of the 2023 Turkish election and to identify, in interpretable terms, where it fails.

%==============================================================
\chapter{Material und Methoden -- Material and Methods}
\label{ch:methods}
%==============================================================

\section{Material}
\label{sec:material}

\subsection{Datensatz -- Dataset}
\label{sec:dataset}

The dataset analysed in this thesis is fully synthetic. It was produced by the simulation pipeline described in Section~\ref{sec:methods}; no human-subject data were collected. The synthetic dataset comprises eight categories of files, all of which are version-controlled and form part of the reproducibility supplement.

The first category is the \textbf{voter-archetype configuration}, stored in the file \path{voter_source_of_truth/synthetic_turkey_simulation.json}. It defines twelve archetypes (Devout Anatolian Loyalist; Secular Urban Professional; Conservative Economically Disillusioned; Alevi-CHP Loyalist; Kurdish Political Voter; Nationalist Grey Wolf; Moderate Nationalist İYİ; Pious Disillusioned Islamist; Young Urban Protest Voter; Earthquake Zone Loyalist; Retired Protest Voter; Cosmopolitan Liberal). For each archetype the file records a research-informed political worldview vector, an emotional baseline vector, a media-diet vector, a behavioural-parameter vector, an initial party-vote distribution, and an initial candidate-vote distribution. The second category is the \textbf{2018 baseline sampling profile} (\path{voter_source_of_truth/2018_baseline_sampling_profile.yaml}), which encodes the party and presidential vote-share targets used to draw a sample that is structurally consistent with the 2018 Turkish elections.

The third category is the \textbf{political-actor configuration}, stored in \path{political_agent/political_agents.yaml}, \path{political_personas.yaml}, \path{credibility_matrix.yaml}, \path{politician_event_responses.yaml}, and \path{movement_state_machine.yaml}. These files define the political actors whose messages reach the voters, the persona-level rhetorical signatures of those actors, the credibility scores that each archetype assigns to each actor, the event-specific response frames, and the time-evolving state of the Kurdish political movement across 2018--2023. The fourth category is the \textbf{event timeline}, stored as \path{events/simulation_ticks.json}. It contains thirty-seven dated event ticks that span the period from the June 2018 election baseline to the May 2023 runoff result reveal.

The fifth category is the generated \textbf{synthetic electorate} itself: three hundred per-agent persona files in \path{souls/agent_001.json} through \path{souls/agent_300.json}. Each file includes the agent's identifier, archetype, demographic identity (age, gender, city, region, education, income, employment), a short biography, a worldview summary, a 2018 baseline memory, a media-diet description, an emotional baseline, behavioural parameters, and the initial candidate probability distributions. The sixth category is the \textbf{episodic memory} archive in the \texttt{db/} directory: three hundred files of the form \path{agent_XXX_episodic_2023.json} hold the agent-specific reflections and memory entries generated during the run, with retrieval counts and source codes; and three hundred companion files of the form \path{agent_XXX_beliefs_2023.json} hold the per-tick belief states.

The seventh category is the main \textbf{analytical output}: \path{outputs/agent_trajectories.csv}. This file contains one row for each (agent, tick) pair, so its length is 300 $\times$ 37 = 11{,}100 rows. The thirty-five columns are: \texttt{tick\_id}, \texttt{sim\_date}, \texttt{agent\_id}, \texttt{archetype\_id}, \texttt{archetype\_name}, \texttt{government\_approval}, \texttt{institutional\_trust}, \texttt{opposition\_trust}, the five affect fields (\texttt{anger}, \texttt{fear}, \texttt{hope}, \texttt{sadness}, \texttt{political\_fatigue}), \texttt{turnout\_probability}, six first-round candidate probabilities (\texttt{first\_round\_Erdogan}, \texttt{first\_round\_Kilicdaroglu}, \texttt{first\_round\_Sinan\_Ogan}, \texttt{first\_round\_Muharrem\_Ince}, \texttt{first\_round\_Other}, \texttt{first\_round\_Undecided}), three runoff probabilities (\texttt{runoff\_Erdogan}, \texttt{runoff\_Kilicdaroglu}, \texttt{runoff\_Abstain\_Invalid\_Undecided}), eight party-preference probabilities (\texttt{party\_AKP}, \texttt{party\_CHP}, \texttt{party\_MHP}, \texttt{party\_IYI}, \texttt{party\_DEM\_HDP\_YSP}, \texttt{party\_YRP}, \texttt{party\_Other}, \texttt{party\_Undecided}), the three argmax shortcut columns (\texttt{first\_round\_top}, \texttt{runoff\_top}, \texttt{party\_top}), and four provenance fields (\texttt{confidence}, \texttt{reason\_codes}, \texttt{reflection}, \texttt{visible\_broadcast\_count}).

The eighth category is the \textbf{run log}: \path{logs/run_2023_20260516_142653_179132.jsonl}, a 18.6 MB JSONL file with 11{,}138 records that contains the full event-by-event trace of the simulation. Auxiliary outputs include \path{outputs/aggregate_candidate_intention.csv}, \path{outputs/aggregate_party_preference.csv}, \path{outputs/reflections.jsonl}, \path{outputs/broadcasts.jsonl}, and the JSON evaluation summary against actual YSK results.

\subsection{Verwendete Tools -- Used Tools}
\label{sec:tools}

The simulation was implemented in Python 3.11. The main library dependencies are \texttt{openai} (for LLM calls to the OpenAI API), \texttt{pyyaml} (for YAML configuration loading), \texttt{pandas} (for data-frame manipulation and CSV export), \texttt{numpy} (for numerical operations), and \texttt{matplotlib} (for figure generation). The LLM provider used for the baseline run was the OpenAI Chat Completions API, with the model identifier \texttt{gpt-4o-mini}, temperature 0.45, and a maximum response budget of 900 tokens per voter-decision call. Parallel voter-decision calls were dispatched through Python's \texttt{concurrent.futures.ThreadPoolExecutor}. Synthetic-population sampling was deterministic, seeded with the integer \texttt{20230528}. Analysis and figure generation were performed with custom Python scripts (\path{scripts/section_d_fresh_audit.py}, \path{scripts/section_k_generate_charts.py}, \path{scripts/tie_handling_sensitivity.py}).

\section{Methoden -- Methods}
\label{sec:methods}

\subsection{Simulation Architecture}
\label{sec:architecture}

Synthetic Turkey is organised as a deterministic orchestration layer wrapped around a stochastic LLM call. Figure~\ref{fig:framework} summarises the four-stage architecture. The first stage constructs the synthetic electorate from configuration files. The second stage initialises each voter as a citizen agent with persistent belief state, episodic memory, and an LLM provider. The third stage iterates the simulation tick by tick: each event tick triggers (a) the generation or retrieval of relevant political broadcasts, (b) the filtering of those broadcasts according to each voter's media diet and credibility matrix, (c) an LLM call that asks the voter to reason about the event, and (d) the deterministic validation and storage of the returned JSON object. The fourth stage performs aggregation and evaluation against actual YSK results.

\begin{figure}[!ht]
\centering
\includegraphics[width=0.82\linewidth]{Fig1_framework.png}
\caption{Synthetic Turkey framework. Deterministic code controls the outer loop; the LLM performs only the voter-level reasoning step.}
\label{fig:framework}
\end{figure}

\subsection{Voter Archetype Construction}
\label{sec:archetype_construction}

The twelve voter archetypes were derived from the Turkish electoral literature reviewed in Chapter~\ref{ch:litreview}. Each archetype is parameterised on a small number of structured vectors. The political-worldview vector encodes positions on conservative--secular, nationalist--cosmopolitan, populist--institutional, authoritarian--libertarian, and pro-EU axes. The behavioural-parameter vector encodes turnout propensity, cross-bloc switching propensity, susceptibility to economic messaging, susceptibility to security messaging, and susceptibility to identity-based mobilization. The emotional baseline encodes resting levels of anger, fear, hope, sadness, and political fatigue. The media-diet vector encodes the share of attention given to pro-government, opposition, and social media. The initial party and candidate distributions are calibrated to the 2018 baseline and reflect the empirical distribution of votes within that archetype's analogue in the literature \citep{carkoglu2024unlikely,aytac2025polarised,baruh2024difficult,sevinc2024afet}.

A central design choice is that the 2018 baseline is treated as memory and identity rather than as a deterministic rule for the 2023 vote. Agents are allowed to revise their preferences over time through LLM reasoning. The sampling frame distributes the three hundred agents across the twelve archetypes such that party and presidential anchors approximate the 2018 YSK distribution: 128 agents anchored to AKP, 68 to CHP, 35 to HDP/DEM, 33 to MHP, 30 to İYİ, and 6 to small parties.

\subsection{Event Timeline and Leakage Prevention}
\label{sec:timeline_methods}

The event timeline contains thirty-seven dated ticks. A crucial methodological feature is the prevention of result leakage at the two election ticks. Both the first-round event and the runoff event are split into a decision phase and a reveal phase. At \path{T030A_first_round_vote_decision} the agent commits its first-round candidate distribution \emph{before} any first-round result is shown; only at \path{T030B_first_round_result_revealed} does the actual result enter the agent's memory. The same applies to \path{T035A_runoff_vote_decision} and \path{T035B_final_result_revealed}. The voter prompt additionally includes a temporal fence: the agent is instructed to reason only from information available up to the simulated date of the current tick, which prevents within-prompt leakage of post-date knowledge.

\subsection{Information Filtering and Political Broadcasts}
\label{sec:broadcasts_methods}

Political communication is represented through controlled broadcast inputs rather than as fully autonomous strategic political agents. For each event tick the system identifies relevant political speakers and event-specific message frames from structured YAML configuration files. The voter does not see the raw broadcast list; instead, each voter sees only a filtered subset that depends on (a) the voter's media diet, (b) the credibility score that the voter's archetype assigns to that political actor in the credibility matrix, (c) the date of the event, and (d) the relevance of the source to the event in question. This produces the realistic property that different voters see different parts of the same event: a pro-government media consumer and a social-media-heavy opposition voter receive different visible messages from the same political moment \citep{karaca2024socialmedia,najafi2024dataset,bozdag2022skeptical}.

\subsection{LLM Voter Reasoning and Prompt Structure}
\label{sec:prompt_structure}

At each tick, each voter agent is presented with a prompt that contains eight ordered blocks. (1) A \emph{temporal fence} instructs the LLM that the agent only knows information available up to the simulated date. (2) A \emph{role instruction} commits the model to first-person voter role-play rather than to third-person political analysis. (3) The agent's \emph{identity and archetype} block names the agent and gives demographic details. (4) The \emph{biography and worldview} block summarises the agent's political identity, media diet, and emotional baseline. (5) The \emph{2018 baseline memory} block records the agent's 2018 vote. (6) The \emph{current state} block lists the agent's previous beliefs, affective state, and recent reflections. (7) The \emph{event and visible broadcasts} block presents the current event and the filtered subset of political messages. (8) The \emph{output schema} block specifies a strict JSON response format.

The output JSON schema requires the following fields with the indicated ranges: \texttt{government\_approval} (1--10), \texttt{institutional\_trust} (1--10), \texttt{opposition\_trust} (1--10), \texttt{anger}, \texttt{fear}, \texttt{hope}, \texttt{sadness}, and \texttt{political\_fatigue} (each 0--1), \texttt{turnout\_probability} (0--1), \texttt{party\_preference} (probability distribution over eight party labels), \texttt{first\_round\_vote\_intention} (probability distribution over six first-round labels), \texttt{runoff\_vote\_intention} (probability distribution over three labels), \texttt{confidence} (low / medium / high), \texttt{reason\_codes} (free-form short tags), and \texttt{reflection} (a short first-person paragraph).

Deterministic code then validates the LLM response. If the response fails JSON parsing or omits required fields, the call is retried up to three times with linear back-off and an instruction that requires strict JSON output. Successfully parsed responses are passed through schema validation: numerical state variables are clamped to their allowed scales, probability distributions are normalised to sum to one over the canonical key sets, and the \texttt{confidence} and \texttt{reason\_codes} fields are sanitised. If all retries fail, the engine can continue with a low-confidence fallback that preserves the agent's previous belief and affect state and marks the case as a provider error. In the baseline run, no fallback was triggered. The model was \texttt{gpt-4o-mini}, the temperature was 0.45, and the maximum response budget was 900 tokens.

\subsection{Aggregation and Scoring Rules}
\label{sec:scoring}

The simulation produces a candidate probability distribution for each agent. These distributions can be aggregated to the population level in two reasonable ways. The first rule, \emph{mean-probability aggregation}, treats each agent's distribution as fractional electoral support:

\begin{equation}
S_c^{\mathrm{mean}} = \frac{1}{N}\sum_{i=1}^{N} p_{ic},
\label{eq:mean}
\end{equation}

where $p_{ic}$ is agent $i$'s probability for candidate $c$. The second rule, \emph{mode-vote aggregation}, assigns each agent to the candidate with the highest probability and counts the discrete choices:

\begin{equation}
v_i = \arg\max_{c} p_{ic},
\label{eq:mode_argmax}
\end{equation}

\begin{equation}
S_c^{\mathrm{mode}} = \frac{1}{N}\sum_{i=1}^{N}\mathbb{1}(v_i=c).
\label{eq:mode_share}
\end{equation}

A non-trivial minority of agents produce tied top probabilities. When two or more candidates share the maximum, a tie-handling rule is required. The implementation's canonical rule resolves ties by a fixed candidate order (Erdoğan first, then Kılıçdaroğlu, then Sinan Oğan, etc.). The thesis additionally reports outcomes under a reverse-order rule (Erdoğan last) and under a Monte Carlo random tie-break with 10{,}000 samples. The divergence between aggregation rules and the sensitivity to tie handling are central to the results.

\subsection{Evaluation Strategy}
\label{sec:evaluation_strategy}

The main external comparison target is the official YSK 2023 result. The first round is compared against Erdoğan 49.52~per cent, Kılıçdaroğlu 44.88~per cent, Sinan Oğan 5.17~per cent, and Muharrem İnce 0.43~per cent. The runoff is compared against Erdoğan 52.18~per cent and Kılıçdaroğlu 47.82~per cent \citep{ysk2023}. Evaluation is conducted at five levels: aggregate vote share, archetype-level voting patterns, demographic breakdowns, longitudinal trust/affect trajectories, and tie-handling sensitivity. This deliberately multi-level evaluation is necessary because a simulation may match the aggregate winner while still failing at the coalition level, or it may fail at the aggregate while still reproducing meaningful bloc-level structure \citep{kang2026dynamok,kreutner2026ep}.

%==============================================================
\chapter{Ergebnisse -- Findings}
\label{ch:findings}
%==============================================================

This chapter reports the empirical findings of the completed baseline run. The chapter is organised descriptively and does not interpret the findings; interpretation is deferred to Chapter~\ref{ch:discussion}. Section~\ref{sec:descriptive} reports descriptive statistics of the synthetic electorate. Section~\ref{sec:aggregate_results} reports aggregate vote shares and the tie-handling sensitivity analysis. Section~\ref{sec:archetypes} reports archetype-level voting patterns. Section~\ref{sec:ogan_results} reports the Sinan Oğan representation gap. Section~\ref{sec:demographics_results} reports demographic breakdowns. Section~\ref{sec:trajectories_results} reports the longitudinal trust and approval trajectories.

\section{Descriptive Statistics}
\label{sec:descriptive}

The synthetic electorate comprises three hundred agents distributed across twelve archetypes. Table~\ref{tab:archetype_counts} reports the per-archetype agent count.

\begin{table}[!ht]
\caption{Synthetic electorate composition: agent count per archetype.}
\label{tab:archetype_counts}
\centering
\begin{tabular}{@{}p{0.55\linewidth}r@{}}
\toprule
Archetype & Agents \\
\midrule
Devout Anatolian Loyalist & 58 \\
Conservative Economically Disillusioned & 33 \\
Secular Urban Professional & 32 \\
Nationalist Grey Wolf (MHP core) & 30 \\
Kurdish Political Voter & 29 \\
Moderate Nationalist (İYİ) & 25 \\
Alevi-CHP Loyalist & 19 \\
Pious Disillusioned Islamist & 18 \\
Retired Protest Voter & 16 \\
Young Urban Protest Voter & 15 \\
Earthquake Zone Loyalist & 14 \\
Cosmopolitan Liberal & 11 \\
\midrule
\textbf{Total} & \textbf{300} \\
\bottomrule
\end{tabular}
\end{table}

The thirty-seven tick timeline spans the period from the June 2018 post-election baseline (ticks T001--T003), through an economic-crisis period (T004--T010), an opposition-coordination phase (T011--T020), a pre-election shock phase (T021--T029), the first-round decision and reveal (T030A--T030B), the runoff campaign (T031--T034), the runoff decision and reveal (T035A--T035B), and a brief post-election period (T036--T037). The complete simulation contains 300 $\times$ 37 = 11{,}100 voter-decision calls.

\section{Aggregate Simulation Results}
\label{sec:aggregate_results}

Figure~\ref{fig:scoring} and Table~\ref{tab:aggregate} report the aggregate vote share under mean-probability and canonical mode-vote aggregation, against the actual YSK outcome.

\begin{figure}[!ht]
\centering
\includegraphics[width=0.82\linewidth]{Fig2_scoring_rules.png}
\caption{First-round (left) and runoff (right) vote shares: actual YSK results against simulated mean-probability and canonical mode-vote shares on the same 300-agent population.}
\label{fig:scoring}
\end{figure}

\begin{table}[!ht]
\caption{Actual and simulated presidential vote shares under mean-probability and canonical mode-vote aggregation.}
\label{tab:aggregate}
\centering
\begin{tabular}{@{}llrrrr@{}}
\toprule
Stage & Source / rule & Erdoğan & Kılıçdaroğlu & Oğan & Other / abstain \\
\midrule
First round & Actual (YSK)              & 49.52 & 44.88 & 5.17 & 0.43 \\
First round & Sim. mean probability     & 41.54 & 49.13 & 1.37 & 7.96 \\
First round & Sim. mode vote            & 54.67 & 45.33 & 0.00 & 0.00 \\
\midrule
Runoff      & Actual (YSK)              & 52.18 & 47.82 & --   & 0.00 \\
Runoff      & Sim. mean probability     & 41.80 & 49.06 & --   & 9.14 \\
Runoff      & Sim. mode vote            & 55.33 & 44.33 & --   & 0.33 \\
\bottomrule
\end{tabular}
\end{table}

Under mean-probability aggregation, the simulation produces a Kılıçdaroğlu plurality in both rounds. Under the canonical mode-vote rule the simulation produces an Erdoğan victory in both rounds. The two aggregation rules therefore yield two different winners from the same agent-level outputs. Table~\ref{tab:tie_sensitivity} reports the sensitivity of the mode-vote result to the tie-handling rule.

\begin{table}[!ht]
\caption{Mode-vote sensitivity to top-probability tie handling. Random tie-break results are based on 10{,}000 Monte Carlo samples.}
\label{tab:tie_sensitivity}
\centering
\footnotesize
\begin{tabular}{@{}llrrrr@{}}
\toprule
Stage & Candidate & Actual (YSK) & Canonical & Reverse & Random mean [95\% CI] \\
\midrule
First round & Erdoğan      & 49.52 & 54.67 & 41.00 & 46.56 [44.67, 48.67] \\
First round & Kılıçdaroğlu & 44.88 & 45.33 & 46.67 & 48.34 [46.33, 50.33] \\
First round & S. Oğan      &  5.17 &  0.00 &  0.33 &  0.39 [0.00, 1.00] \\
First round & M. İnce      &  0.43 &  0.00 &  0.00 &  0.00 [0.00, 0.00] \\
First round & Other        &  n/a &  0.00 &  0.33 &  0.17 [0.00, 0.33] \\
First round & Undecided    &  n/a &  0.00 & 11.67 &  4.55 [2.67, 6.33] \\
\midrule
Runoff      & Erdoğan      & 52.18 & 55.33 & 40.33 & 46.45 [44.33, 48.67] \\
Runoff      & Kılıçdaroğlu & 47.82 & 44.33 & 43.33 & 46.60 [44.67, 48.67] \\
Runoff      & Abstain/Inv./Undec. & n/a & 0.33 & 16.33 &  6.95 [4.67, 9.33] \\
\bottomrule
\end{tabular}
\end{table}

Fifty of three hundred first-round agents (16.7~per cent) and fifty-four of three hundred runoff agents (18.0~per cent) produce tied top probabilities at the decision tick. The canonical mode-vote rule resolves these ties in Erdoğan's favour. Under the reverse-order rule and under random tie-breaking, the first-round mode-vote winner shifts to Kılıçdaroğlu. Under random tie-breaking the runoff is nearly indistinguishable between the two candidates.

\section{Archetype-Level Voting Patterns}
\label{sec:archetypes}

Figure~\ref{fig:archetype_breakdown} reports the first-round mode-vote breakdown by archetype.

\begin{figure}[!ht]
\centering
\includegraphics[width=0.82\linewidth]{Fig3_archetype_mode_vote.png}
\caption{First-round mode-vote by archetype: share of agents in each archetype whose top-ranked first-round candidate was the listed candidate.}
\label{fig:archetype_breakdown}
\end{figure}

The Devout Anatolian Loyalist, Nationalist Grey Wolf, and Earthquake Zone Loyalist archetypes vote entirely for Erdoğan. The Secular Urban Professional, Alevi-CHP Loyalist, Young Urban Protest Voter, and Cosmopolitan Liberal archetypes vote entirely for Kılıçdaroğlu. The Kurdish Political Voter archetype also consolidates behind Kılıçdaroğlu. Conservative Economically Disillusioned and Pious Disillusioned Islamist voters remain predominantly Erdoğan-oriented but show some erosion. Moderate Nationalist İYİ voters lean strongly toward Kılıçdaroğlu while retaining a pro-Erdoğan minority. Retired Protest Voters split evenly between Erdoğan and Kılıçdaroğlu.

\section{The Oğan Representation Gap}
\label{sec:ogan_results}

The simulation underestimates Sinan Oğan's first-round support. The actual Oğan vote share was 5.17~per cent; the simulated mean-probability share is 1.37~per cent and the simulated mode-vote share is zero. In a 300-agent population, 5.17~per cent corresponds to approximately fifteen agents. The simulation produces no archetype-level bloc for which Oğan is the top candidate. The tie analysis shows that Oğan appears in only a small number of tied top-probability distributions (two first-round three-way ties involving Erdoğan, Oğan, and Undecided, and one two-way tie involving Erdoğan and Oğan), but never as the unique most-probable candidate. The Sinan Oğan miss is therefore a representation gap rather than a marginal numerical error: the simulation contains nationalist archetypes (Nationalist Grey Wolf, Moderate Nationalist İYİ) but does not contain a dedicated nationalist-protest pathway.

\section{Demographic Breakdowns}
\label{sec:demographics_results}

Figure~\ref{fig:demographics} reports demographic breakdowns of first-round mode-vote by region, education, age bracket, and city.

\begin{figure}[!ht]
\centering
\includegraphics[width=0.82\linewidth]{Fig5_demographic_breakdowns.png}
\caption{Demographic breakdown of first-round mode-vote: by region (top left), education (top right), age bracket (bottom left), and city (bottom right, top ten by N).}
\label{fig:demographics}
\end{figure}

Marmara, Aegean, Istanbul, Ankara, and Izmir lean toward Kılıçdaroğlu. Central Anatolia, the Black Sea, Kayseri, Konya, Sivas, Erzurum, Yozgat, and Mersin lean toward Erdoğan. Younger voters lean strongly toward Kılıçdaroğlu while older groups are more Erdoğan-oriented. Graduate and university-student categories are fully Kılıçdaroğlu in the simulation; primary, high-school, and vocational groups lean more Erdoğan. Several groups are reported as fully one candidate in the synthetic data and some city-level groups contain few agents.

\section{Longitudinal State Trajectories}
\label{sec:trajectories_results}

Figure~\ref{fig:trajectories} reports mean government approval, institutional trust, and opposition trust across the thirty-seven simulation ticks.

\begin{figure}[!ht]
\centering
\includegraphics[width=0.82\linewidth]{Fig4_trust_trajectories.png}
\caption{Mean government approval, institutional trust, and opposition trust across 37 simulation ticks. Background bands mark major event clusters for orientation.}
\label{fig:trajectories}
\end{figure}

Government approval and institutional trust decline gradually across the timeline, with a more pronounced decline around the earthquake cluster (T021--T029). Opposition trust rises modestly across the same period. The magnitudes of the trajectory shifts are smaller than one might expect on the basis of the qualitative content of voter reflections.

%==============================================================
\chapter{Diskussion und Schlussfolgerung -- Discussion and Conclusion}
\label{ch:discussion}
%==============================================================

This chapter interprets the findings of Chapter~\ref{ch:findings} in relation to the literature reviewed in Chapter~\ref{ch:litreview}.

\section{Aggregation Sensitivity as a Validity Problem}
\label{sec:disc_aggregation}

The central methodological finding of this thesis is that the same agent-level output supports two different election interpretations depending on how individual probability distributions are converted into a population-level vote share. Under mean-probability aggregation the model produces a Kılıçdaroğlu plurality; under canonical mode-vote aggregation the model produces an Erdoğan plurality; under random tie-handling the mode-vote result drifts back toward Kılıçdaroğlu. This is consistent with the warning in \citet{kang2026dynamok} that apparent winner accuracy in an LLM-agent electoral simulation can mask substantial methodological fragility, and with the broader concern raised in \citet{qu2024performance} and \citet{rettenberger2025politicalbias} that LLM-generated probabilities should not be treated as calibrated electoral probabilities without external validation.

In substantive terms, the canonical mode-vote rule reproduces the correct first-round winner. However, that reproduction is partly an artefact of the deterministic tie-handling order. Approximately one in six first-round agents produced a tied top probability at the decision tick, and the canonical implementation systematically resolves those ties in Erdoğan's favour. Reversing the order or randomising the tie-break shifts the simulated winner. The implication is that LLM electoral simulations should report \emph{both} the aggregation rule and the tie-handling rule explicitly, and should treat the divergence between rules as substantive information about the calibration of the underlying probabilities \citep{kreutner2026ep}.

\section{Coalition-Level Plausibility}
\label{sec:disc_coalition}

At the archetype level, the simulation reproduces several patterns that are well established in the Turkish empirical literature. The Devout Anatolian Loyalist archetype consolidates behind Erdoğan, consistent with the finding of \citet{carkoglu2024unlikely} that Erdoğan's 2023 support was concentrated among right-of-centre and highly religious voters and resilient against negative economic evaluations. The Kurdish Political Voter archetype consolidates behind Kılıçdaroğlu, consistent with the finding of \citet{baruh2024difficult} that Kılıçdaroğlu had a comparative advantage among HDP-affinity voters relative to nationalist-coded alternatives, and with the institutional account in \citet{ilbiz2023securitization} of how securitisation pressure shaped opposition-side coordination.

Particularly important is the Earthquake Zone Loyalist archetype, which the simulation correctly preserves on the Erdoğan side. A naïve retrospective-accountability model would have predicted opposition switching among disaster-affected voters. The simulation instead reproduces the more nuanced finding of \citet{sevinc2024afet}, who argue that earthquake-affected voters engaged in mediated retrospective voting, weighing services, aid, and political identity alongside disaster response. The simulation therefore captures a real empirical puzzle of the 2023 election: that disaster exposure did not automatically translate into incumbent punishment.

\section{Third-Party Collapse and Over-Homogenisation}
\label{sec:disc_thirdparty}

The simulation's clearest failure is the absence of a discrete Oğan-supporting bloc, which directly corresponds to the third-party salience-collapse pattern reported by \citet{kang2026dynamok} for the Korean case. Both simulations contain nationalist archetypes, but neither contains a sufficiently distinct nationalist-protest pathway to produce a discrete third candidate at the mode-vote level. In the Turkish case, this gap is consistent with the empirical work of \citet{gerim2025victory} on anti-immigration mobilisation and the Victory Party, and with \citet{karaca2024socialmedia} on the role of social-media mobilisation, both of which suggest that the Oğan/ATA pathway depends on a specific protest-vote audience that cannot be reduced to a residual probability inside the Grey Wolf or İYİ archetypes.

A related but distinct problem is the over-homogenisation of certain opposition archetypes. The Secular Urban Professional, Alevi-CHP Loyalist, Young Urban Protest Voter, and Cosmopolitan Liberal archetypes vote entirely for Kılıçdaroğlu in the mode-vote result, which is unrealistic at the individual level even though it is directionally correct at the coalition level. This over-homogenisation likely reflects two combined design choices: the strong prior tie of these archetypes to opposition voting in the 2018 baseline, and the absence of agent-internal disagreement that would, in a real population, produce a tail of cross-pressured or abstaining voters within the same archetype. The pattern is consistent with the broader concern in \citet{aytac2025polarised} that affective polarisation tends to lock voter blocs in place, but the simulation's representation of it is too sharp.

The trust and affect trajectories show a separate diagnostic concern: the damping problem. Government approval and institutional trust decline in the correct direction around the earthquake cluster, but the magnitudes of these shifts are smaller than the qualitative content of the voter reflections would suggest. This is consistent with the broader observation in \citet{ziems2024css} and \citet{qu2024performance} that LLM outputs in structured numerical fields can be less responsive than LLM outputs in free-text fields. A model may therefore appear politically responsive in natural-language reflection while remaining numerically inert in its state variables.

\section{Limitations}
\label{sec:limitations}

This study has several important limitations. First, the completed simulation is a single stochastic realisation at non-zero temperature. The tie-handling Monte Carlo analysis estimates uncertainty only for tied top-probability conversion inside the completed run; it does not estimate LLM run-to-run variance. Repeated full simulations across model versions, temperatures, prompt variants, and population seeds would be required for stronger robustness claims. Second, the primary mode-vote result is sensitive to tie handling, as discussed above. Third, the synthetic electorate is not generated from licensed survey microdata; the voter archetypes are research-informed but synthetic. Fourth, latent LLM knowledge of the 2023 Turkish election cannot be fully removed by temporal fencing alone, even though the leakage-prevention design substantially reduces the risk. Fifth, researcher priors are embedded in archetype definitions, sampling weights, credibility matrices, event descriptions, broadcast frames, and prompt wording. Sixth, political broadcasts are controlled inputs rather than fully autonomous agents. Seventh, the model lacks a dedicated nationalist-protest pathway for Sinan Oğan. Eighth, causal event attribution is limited by the absence of an ablation design. Ninth, numerical state trajectories are damped relative to qualitative reflection. Tenth, subgroup analysis is limited by sample size: percentages for small archetype $\times$ city subgroups should be read cautiously.

\section{Conclusion and Future Work}
\label{sec:conclusion}

This thesis has presented Synthetic Turkey, a source-grounded LLM-agent simulation of the 2023 Turkish presidential election, and has used the simulation as a diagnostic instrument rather than as a predictive tool. The simulation reproduces several substantively important patterns of the 2023 election at the coalition level, including the persistence of Erdoğan loyalty in the earthquake zone, the Kurdish strategic consolidation behind Kılıçdaroğlu, and the urban--rural cleavage. It also exposes several methodologically important failure modes: aggregation sensitivity, tie-handling sensitivity, third-party representation collapse, over-homogenisation of opposition archetypes, and damped numerical state trajectories.

The broader implication of these findings is that LLM-agent electoral simulations should be evaluated as diagnostic instruments before they are treated as predictive models. Their value lies not only in whether they match an election result, but in where they fail, which assumptions create those failures, and how those failures reveal the structure of the model. Future work should pursue four concrete extensions. First, the simulation should be repeated across multiple LLM models, temperatures, and prompt variants so that LLM run-to-run variance can be quantified directly, addressing the central limitation of the present study. Second, a dedicated nationalist-protest archetype should be added and calibrated to the Sinan Oğan / ATA Alliance / Victory Party empirical literature to test whether the third-party collapse pattern is structural or an artefact of the present archetype set. Third, the broadcast layer should be extended toward autonomous but constrained LLM political agents, while preserving source control and leakage prevention. Fourth, the within-archetype heterogeneity should be increased so that the over-homogenisation pattern can be tested against more empirically realistic intra-bloc dispersion. The combination of these four extensions would substantially strengthen the diagnostic value of an already promising methodology.

%==============================================================
% References
%==============================================================
\chapter{Literaturverzeichnis -- References}
\label{ch:references}

\begin{thebibliography}{99}

\bibitem[Aher et~al.(2023)]{aher2023simulating}
Aher, G.~V., Arriaga, R.~I., \& Kalai, A.~T. (2023). Using large language models to simulate multiple humans and replicate human subject studies. In \emph{Proceedings of the 40th International Conference on Machine Learning} (PMLR, Vol.~202).

\bibitem[Argyle et~al.(2023)]{argyle2023out}
Argyle, L.~P., Busby, E.~C., Fulda, N., Gubler, J.~R., Rytting, C., \& Wingate, D. (2023). Out of one, many: Using language models to simulate human samples. \emph{Political Analysis, 31}(3), 337--351. \url{https://doi.org/10.1017/pan.2023.2}

\bibitem[Aytaç et~al.(2025)]{aytac2025polarised}
Aytaç, S.~E., Irmak, S., \& Söylemez-Karakoç, B. (2025). The polarised electorate: Affective polarisation and electoral loyalty in Turkey's 2023 critical election. \emph{South European Society and Politics}, advance online publication. \url{https://doi.org/10.1080/13608746.2025.2597287}

\bibitem[Baruh and Çarkoğlu(2024)]{baruh2024difficult}
Baruh, L., \& Çarkoğlu, A. (2024). Difficult choices: Choosing the candidate of the Nation Alliance in the 2023 Turkish presidential election. \emph{Southeast European and Black Sea Studies, 24}(2), 319--341. \url{https://doi.org/10.1080/14683857.2024.2324523}

\bibitem[Bozdağ and Koçer(2022)]{bozdag2022skeptical}
Bozdağ, Ç., \& Koçer, S. (2022). Skeptical inertia in the face of polarization: News consumption and misinformation in Turkey. \emph{Media and Communication, 10}(2), 222--232. \url{https://doi.org/10.17645/mac.v10i2.5057}

\bibitem[Çarkoğlu(2024)]{carkoglu2024unlikely}
Çarkoğlu, A. (2024). The unlikely survival of Erdoğan in Turkey's May 2023 elections. \emph{JCMS: Journal of Common Market Studies, 62}(S1), 162--173. \url{https://doi.org/10.1111/jcms.13655}

\bibitem[Çınkara and Çoşkun(2025)]{cinkara2025erdogan}
Çınkara, G., \& Çoşkun, B. (2025). Erdoğan's electoral success: Geopolitics, social class and new nationalism. \emph{Liberal Düşünce, 30}(117), 45--72. \url{https://doi.org/10.36484/liberal.1527011}

\bibitem[Demir(2024)]{demir2024iyi}
Demir, O. (2024). Siyasal parti teorileri bağlamında bir değerlendirme: İYİ Parti. \emph{Uluslararası Yönetim Akademisi Dergisi, 7}(4), 1148--1162. \url{https://doi.org/10.33712/mana.1544435}

\bibitem[Epstein(2006)]{epstein2006generative}
Epstein, J.~M. (2006). \emph{Generative social science: Studies in agent-based computational modeling}. Princeton University Press.

\bibitem[Gao et~al.(2024)]{gao2024llmabm}
Gao, C., Lan, X., Li, N., Yuan, Y., Ding, J., Zhou, Z., Xu, F., \& Li, Y. (2024). Large language models empowered agent-based modeling and simulation: A survey and perspectives. \emph{Humanities and Social Sciences Communications, 11}, 1259. \url{https://doi.org/10.1057/s41599-024-03611-3}

\bibitem[Gerim(2025)]{gerim2025victory}
Gerim, G. (2025). Influence of anti-immigration attitudes on voter behavior and its limits: Turkish politics and the case of the Victory Party. \emph{Marmara Üniversitesi Siyasal Bilimler Dergisi, 13}(1), 181--199. \url{https://doi.org/10.14782/marmarasbd.1574198}

\bibitem[Hafızoğlu(2021)]{hafizoglu2021z}
Hafızoğlu, Y. (2021). Z kuşağı seçmen davranışları. \emph{Eğitim ve Yeni Yaklaşımlar Dergisi, 4}(2), 141--164. \url{https://doi.org/10.52974/jena.1035189}

\bibitem[Hwang et~al.(2023)]{hwang2023aligning}
Hwang, E., Majumder, B.~P., \& Tandon, N. (2023). Aligning language models to user opinions. In \emph{Findings of the Association for Computational Linguistics: EMNLP 2023} (pp.~5906--5919).

\bibitem[Ilbiz and Kaunert(2023)]{ilbiz2023securitization}
Ilbiz, E., \& Kaunert, C. (2023). Securitization, fear politics, and the formation of an opposition alliance in competitive authoritarian regimes. \emph{Democratization, 30}(8), 1604--1624. \url{https://doi.org/10.1080/13510347.2023.2252346}

\bibitem[Kang(2026)]{kang2026dynamok}
Kang, S. (2026). \emph{Diagnosing Korean-language LLM political bias via census-grounded agent simulation} (arXiv:2605.18395) [Preprint]. arXiv.

\bibitem[Karaca(2024)]{karaca2024socialmedia}
Karaca, M.~F. (2024). 2023 Cumhurbaşkanlığı seçiminde sosyal medya kullanımı: Instagram ve X üzerine bir inceleme. \emph{İnönü Üniversitesi Uluslararası Sosyal Bilimler Dergisi, 13}(1), 136--175. \url{https://doi.org/10.54282/inijoss.1454160}

\bibitem[Kreutner et~al.(2026)]{kreutner2026ep}
Kreutner, M., Lutz, M., \& Strohmaier, M. (2026). Persona-driven simulation of voting behavior in the European Parliament with large language models. In \emph{Findings of the Association for Computational Linguistics: EACL 2026}.

\bibitem[Kuru(2017)]{kuru2017religious}
Kuru, A.~T. (2017). Constitution-writing and religious divisions in Turkey. In A.~Shachar, R.~Hirschl, \& I.~Shapiro (Eds.), \emph{Constitution writing, religion and democracy}. Cambridge University Press.

\bibitem[Lu et~al.(2024)]{lu2024gabs}
Lu, Y., Aleta, A., Du, C., Shi, L., \& Moreno, Y. (2024). LLMs and generative agent-based models for complex systems research. \emph{Physics of Life Reviews, 51}, 283--293. \url{https://doi.org/10.1016/j.plrev.2024.10.013}

\bibitem[Macal and North(2010)]{macal2010tutorial}
Macal, C.~M., \& North, M.~J. (2010). Tutorial on agent-based modelling and simulation. \emph{Journal of Simulation, 4}(3), 151--162. \url{https://doi.org/10.1057/jos.2010.3}

\bibitem[Najafi et~al.(2024)]{najafi2024dataset}
Najafi, A., Mugurtay, N., Zouzou, Y., Demirci, E., Demirkiran, S., Karadeniz, H.~A., \& Varol, O. (2024). First public dataset to study 2023 Turkish general election. \emph{Scientific Reports, 14}, 8794. \url{https://doi.org/10.1038/s41598-024-58006-w}

\bibitem[Özbudun(2013)]{ozbudun2013cleavages}
Özbudun, E. (2013). \emph{Party politics and social cleavages in Turkey}. Lynne Rienner Publishers.

\bibitem[Park et~al.(2023)]{park2023generative}
Park, J.~S., O'Brien, J.~C., Cai, C.~J., Morris, M.~R., Liang, P., \& Bernstein, M.~S. (2023). Generative agents: Interactive simulacra of human behavior. In \emph{Proceedings of the 36th Annual ACM Symposium on User Interface Software and Technology}. \url{https://doi.org/10.1145/3586183.3606763}

\bibitem[Qu and Wang(2024)]{qu2024performance}
Qu, Y., \& Wang, J. (2024). Performance and biases of large language models in public opinion simulation. \emph{Humanities and Social Sciences Communications, 11}, 1095. \url{https://doi.org/10.1057/s41599-024-03609-x}

\bibitem[Rettenberger et~al.(2025)]{rettenberger2025politicalbias}
Rettenberger, L., Reischl, M., \& Schutera, M. (2025). Assessing political bias in large language models. \emph{Journal of Computational Social Science, 8}, Article~42. \url{https://doi.org/10.1007/s42001-025-00376-w}

\bibitem[Sevinç and Göz(2024)]{sevinc2024afet}
Sevinç, H., \& Göz, M.~M. (2024). Doğal afetlerin seçmen davranışı üzerine etkisi: 6 Şubat Kahramanmaraş depremi örneği. \emph{Akademik Yaklaşımlar Dergisi, 15}(1), 790--821. \url{https://doi.org/10.54688/ayd.1411908}

\bibitem[Supreme Election Council(2023)]{ysk2023}
Supreme Election Council of Türkiye (YSK). (2023). \emph{Official results of the 2023 Turkish presidential election}. YSK Official Election Data. Accessed 23 May 2026.

\bibitem[Yeğen(2022)]{yegen2022kurdish}
Yeğen, M. (2022). \emph{Erdoğan and the Turkish opposition revisit the Kurdish question} (SWP Comment). Stiftung Wissenschaft und Politik.

\bibitem[Yıldırım(2025)]{yildirim2025sege}
Yıldırım, U.~T. (2025). Socio-economic development and voter preferences in Türkiye: The 28 May 2023 presidential election. \emph{Akademik Araştırmalar ve Çalışmalar Dergisi, 17}(32), 165--180. \url{https://doi.org/10.20990/kilisiibfakademik.1653695}

\bibitem[Zeynel and Savut(2024)]{zeynel2024duygular}
Zeynel, Ş.~Ş., \& Savut, E. (2024). Politik duyguların kültür ve parti eğilimleri ile ilişkisi: 2023 genel seçim kampanyası üzerinden nitel bir analiz. \emph{KMÜ Sosyal ve Ekonomik Araştırmalar Dergisi, 26}(46), 502--522.

\bibitem[Zhang et~al.(2024)]{zhang2024electionsim}
Zhang, X., Lin, J., Sun, L., Qi, W., Yang, Y., Chen, Y., Lyu, H., Mou, X., Chen, S., Luo, J., Huang, X., Tang, S., \& Wei, Z. (2024). \emph{ElectionSim: Massive population election simulation powered by large language model driven agents} (arXiv:2410.20746) [Preprint]. arXiv.

\bibitem[Zhou et~al.(2025)]{zhou2025flockvote}
Zhou, L., Xu, Y., Wang, Z., \& Wang, D. (2025). \emph{FlockVote: LLM-empowered agent-based modeling for simulating U.S. presidential elections} [Preprint]. arXiv.

\bibitem[Ziems et~al.(2024)]{ziems2024css}
Ziems, C., Held, W., Shaikh, O., Chen, J., Zhang, Z., \& Yang, D. (2024). Can large language models transform computational social science? \emph{Computational Linguistics, 50}(1), 237--291. \url{https://doi.org/10.1162/coli_a_00502}

\end{thebibliography}

%==============================================================
% Appendices
%==============================================================
\appendix
\renewcommand{\thechapter}{\Alph{chapter}}

\chapter{Voter Archetype Profiles}
\label{app:archetype_profiles}

Table~\ref{tab:archetype_profiles} summarises the twelve voter archetypes used in the simulation. Each archetype is characterised by a short profile description, its dominant political orientation, and the literature anchors that justify its design.

\begin{landscape}
\small
\begin{longtable}{@{}p{3.5cm}p{1.5cm}p{10cm}p{4.5cm}@{}}
\caption{Voter archetype profiles used in the Synthetic Turkey simulation.}\label{tab:archetype_profiles}\\
\toprule
\textbf{Archetype} & \textbf{N} & \textbf{Profile} & \textbf{Key literature anchors} \\
\midrule
\endfirsthead
\multicolumn{4}{@{}l}{\textit{(continued from previous page)}}\\
\toprule
\textbf{Archetype} & \textbf{N} & \textbf{Profile} & \textbf{Key literature anchors} \\
\midrule
\endhead
Devout Anatolian Loyalist & 58 & Religious, conservative voter strongly attached to the AKP coalition; resilient against economic dissatisfaction. & \citet{carkoglu2024unlikely,aytac2025polarised} \\
Conservative Economically Disillusioned & 33 & Right-of-centre voter under acute economic pressure; cross-pressured between cultural identity and material grievance. & \citet{carkoglu2024unlikely,yildirim2025sege,cinkara2025erdogan} \\
Secular Urban Professional & 32 & Highly educated urban voter aligned with the secular CHP tradition. & \citet{aytac2025polarised,carkoglu2024unlikely} \\
Nationalist Grey Wolf (MHP core) & 30 & Hard-nationalist voter anchored to the People's Alliance; salience on security and anti-immigration. & \citet{carkoglu2024unlikely,gerim2025victory,karaca2024socialmedia} \\
Kurdish Political Voter & 29 & HDP/YSP-affinity voter; engaged in strategic anti-incumbent coordination behind Kılıçdaroğlu. & \citet{baruh2024difficult,ilbiz2023securitization,yegen2022kurdish} \\
Moderate Nationalist (İYİ) & 25 & Centre-right nationalist voter uncomfortable with Kılıçdaroğlu but inside the Nation Alliance. & \citet{baruh2024difficult,demir2024iyi} \\
Alevi-CHP Loyalist & 19 & Alevi-identifying voter with strong secular orientation and negative predisposition toward Erdoğan. & \citet{carkoglu2024unlikely,kuru2017religious,ozbudun2013cleavages} \\
Pious Disillusioned Islamist & 18 & Religious voter dissatisfied with AKP rule; partially open to YRP and small-party defection. & \citet{carkoglu2024unlikely,cinkara2025erdogan} \\
Retired Protest Voter & 16 & Fixed-income elderly voter with strong economic grievance; treated as exploratory pathway. & \citet{carkoglu2024unlikely,yildirim2025sege} \\
Young Urban Protest Voter & 15 & Generation-Z voter with high social-media exposure and lower traditional party loyalty. & \citet{hafizoglu2021z,karaca2024socialmedia,aytac2025polarised} \\
Earthquake Zone Loyalist & 14 & Voter in the eleven earthquake-affected provinces who engages in mediated retrospective voting. & \citet{sevinc2024afet,carkoglu2024unlikely} \\
Cosmopolitan Liberal & 11 & Highly educated, internationally connected secular voter; strongly opposition-leaning. & \citet{aytac2025polarised,ozbudun2013cleavages} \\
\bottomrule
\end{longtable}
\normalsize
\end{landscape}

\chapter{Prompt Template and JSON Output Schema}
\label{app:prompt_schema}

This appendix reproduces the exact prompt that the engine sends to the language model on every voter-decision call, together with the JSON output schema that the model is required to return. The text below is the literal template assembled by the \path{agents/citizen_agent.py} module of the simulation codebase. Placeholders in curly braces are substituted at runtime with values from the agent's persona, the current event, the visible political broadcasts, and the agent's running belief and affect state. Sections labelled \texttt{ROLE}, \texttt{IDENTITY}, \texttt{PERSONA}, and so on appear in the prompt in the order shown, separated by blank lines.

\section*{B.1\quad The full voter-decision prompt}

\begin{lstlisting}[basicstyle=\ttfamily\scriptsize,breaklines=true,frame=single,columns=fullflexible,xleftmargin=0.5em]
You are this synthetic voter. You only know information available up to
the current simulated date: {current_date}. Do not use future knowledge.
Stay consistent with your persona, but you may gradually change your
views if events affect you.

ROLE
Role-play as the synthetic voter below, not as a political analyst.
Reason from their identity, social context, media diet, and memory.

IDENTITY
agent_id: {agent_id}
archetype: {archetype_id} - {archetype_name}
age: {age}
gender: {gender}
city/region: {city} / {region}
education: {education_level}
income: {income_bracket}
employment: {employment_status}

PERSONA
biography: {short_biography}
political identity: {political_identity_summary}
worldview summary: {worldview_summary}
media diet summary: {media_diet_summary}
social context summary: {social_context_summary}

2018 BASELINE MEMORY
{baseline_2018_memory_summary}

NUMERIC GROUNDING
political_worldview: {political_worldview_json}
media_diet: {media_diet_json}
behavioral_variables: {behavioral_variables_json}

CURRENT STATE
{beliefs.to_context_string()}
{affect.to_context_string()}

SHORT MEMORY
{top_5_retrieved_episodic_memories}

CURRENT EVENT
{
  "tick_id": "{tick_id}",
  "date": "{date}",
  "title": "{title}",
  "category": "{category}",
  "summary": "{summary}",
  "notes_for_llm_agents": "{notes_for_llm_agents}"
}

VISIBLE POLITICAL BROADCASTS
{
  "broadcasts": [
    { "agent_name": "{actor_name}",
      "party_or_movement": "{actor_party}",
      "paraphrased_message": "{message_text}",
      "tone": ["{tone_tag_1}", "{tone_tag_2}"],
      "credibility_score": {credibility_0_to_1},
      "likely_effect": "{persuade|mobilize_base|reinforce|backlash|ignored}",
      "source_agent_id": "{actor_id}"
    }
    ... (up to 5 broadcasts, filtered by credibility * media exposure)
  ]
}

VISIBLE PEER SOCIAL CONTEXT
- {peer_reflection_1}
- {peer_reflection_2}
- {peer_reflection_3}
- {peer_reflection_4}
- {peer_reflection_5}

TASK
Return strict JSON only. No markdown, no commentary.
decision_kind: {decision_kind}

Required schema:
{
  "reflection":          "1-3 sentence first-person reflection",
  "government_approval": <number 1-10>,
  "institutional_trust": <number 1-10>,
  "opposition_trust":    <number 1-10>,
  "anger":               <number 0-1>,
  "fear":                <number 0-1>,
  "hope":                <number 0-1>,
  "sadness":             <number 0-1>,
  "political_fatigue":   <number 0-1>,
  "party_preference": {
    "AKP": <0-1>, "CHP": <0-1>, "MHP": <0-1>, "IYI": <0-1>,
    "DEM_HDP_YSP": <0-1>, "YRP": <0-1>,
    "Other": <0-1>, "Undecided": <0-1>
  },
  "first_round_vote_intention": {
    "Erdogan": <0-1>, "Kilicdaroglu": <0-1>,
    "Sinan_Ogan": <0-1>, "Muharrem_Ince": <0-1>,
    "Other": <0-1>, "Undecided": <0-1>
  },
  "runoff_vote_intention": {
    "Erdogan": <0-1>, "Kilicdaroglu": <0-1>,
    "Abstain_Invalid_Undecided": <0-1>
  },
  "turnout_probability": <number 0-1>,
  "reason_codes":        ["short_code"],
  "confidence":          "low|medium|high"
}
\end{lstlisting}

\section*{B.2\quad What each placeholder receives at runtime}

The placeholders shown in curly braces above are populated deterministically from configuration files immediately before the call. Their sources are listed in Table~\ref{tab:placeholder_sources}.

\begin{table}[!htbp]
\caption{Sources of the runtime placeholders in the voter prompt.}
\label{tab:placeholder_sources}
\centering
\footnotesize
\begin{tabular}{@{}p{0.34\linewidth}p{0.58\linewidth}@{}}
\toprule
\textbf{Placeholder} & \textbf{Source} \\
\midrule
\texttt{\{current\_date\}} & The simulated date attached to the current tick. \\
\texttt{\{agent\_id\}, \{archetype\_id\}, \{archetype\_name\}} & From the soul file \path{souls/agent_NNN.json}. \\
\texttt{\{age\}, \{gender\}, \{city\}, \{region\}, \{education\_level\}, \{income\_bracket\}, \{employment\_status\}} & From the \texttt{identity} block of the soul. \\
\texttt{\{short\_biography\}, \{political\_identity\_summary\}, \{worldview\_summary\}, \{media\_diet\_summary\}, \{social\_context\_summary\}, \{baseline\_2018\_memory\_summary\}} & From the \texttt{persona} block of the soul. \\
\texttt{\{political\_worldview\_json\}, \{media\_diet\_json\}, \{behavioral\_variables\_json\}} & From the \texttt{numeric\_profile} block of the soul, serialised as JSON. \\
\texttt{\{tick\_id\}, \{date\}, \{title\}, \{category\}, \{summary\}, \{notes\_for\_llm\_agents\}} & From the event timeline file \path{events/simulation_ticks.json}. \\
\texttt{visible broadcasts} & Generated by the political-broadcast agent, then filtered by the agent's archetype credibility scores and media diet (top 5 retained per tick). \\
\texttt{peer reflections} & The most recent reflections from up to 5 randomly-sampled other agents, drawn from the previous tick's stored reflection log. \\
\texttt{\{beliefs.to\_context\_string()\}, \{affect.to\_context\_string()\}} & Serialised forms of the agent's most recent belief and affect state, computed deterministically from the prior tick's validated JSON response. \\
\texttt{\{decision\_kind\}} & One of \texttt{normal}, \texttt{first\_round\_vote}, or \texttt{runoff\_vote}, set by the engine based on the current tick. \\
\bottomrule
\end{tabular}
\end{table}

\section*{B.3\quad Election-safe variants of the prompt}

At the two election decision ticks the prompt is materially different. At \path{T030A_first_round_vote_decision} and \path{T035A_runoff_vote_decision}, the engine substitutes a stripped event payload in which the \texttt{summary} field reads \emph{``Synthetic voter decision tick. Actual election result is hidden from agents.''} and the fields \texttt{affected\_dimensions}, \texttt{emotional\_impact}, \texttt{candidate\_effect\_hint}, and \texttt{affected\_archetypes} are emptied. The agent thus produces its first-round and runoff distributions without ever having seen the actual YSK result. Only at the matching reveal ticks (\path{T030B_first_round_result_revealed} and \path{T035B_final_result_revealed}) is the full event payload, including the official 2023 percentages, restored.

\section*{B.4\quad Validation applied to the LLM response}

Every JSON object returned by the LLM is processed by a deterministic validator before storage. The validator (a) checks that the fifteen required fields are present, (b) clamps numerical values to their allowed scales, (c) normalises the three probability distributions to sum to one over the canonical key sets, (d) sanitises the \texttt{confidence} field to one of \texttt{low}, \texttt{medium}, \texttt{high}, and (e) wraps \texttt{reason\_codes} in a list if a bare string was returned. Responses that cannot be parsed as JSON or that omit required fields trigger up to three retries with linear back-off and a clarifying instruction; in the baseline run reported here, all 11{,}100 voter-decision calls were resolved within these retries and no low-confidence fallback was triggered.

\chapter{Simulation Event Timeline}
\label{app:event_timeline}

Table~\ref{tab:event_phases} summarises the seven phases of the thirty-seven-tick simulation event timeline.

\begin{table}[!ht]
\caption{Summary of simulation event phases (37 ticks total).}
\label{tab:event_phases}
\centering
\footnotesize
\begin{tabular}{@{}p{0.24\linewidth}p{0.48\linewidth}p{0.20\linewidth}@{}}
\toprule
Phase & Content & Tick range \\
\midrule
2018 baseline           & Post-2018 electoral memory and initial political alignment & T001--T003 \\
Economic crisis period  & Currency crisis, inflation, economic dissatisfaction, government performance updates & T004--T010 \\
Opposition coordination & 2019 local elections, alliance formation, candidate debates, Kurdish strategic context & T011--T020 \\
Pre-election shocks     & Earthquake, campaign dynamics, polling shifts, nationalist fragmentation & T021--T029 \\
First round             & First-round decision tick followed by first-round result reveal & T030A--T030B \\
Runoff campaign         & Oğan endorsement, opposition nationalist turn, second-round mobilisation & T031--T034 \\
Runoff                  & Runoff decision tick followed by final result reveal & T035A--T035B \\
Post-election           & Final memory update and result interpretation & T036--T037 \\
\bottomrule
\end{tabular}
\end{table}

The two election ticks are split into decision and reveal phases. Agents make their first-round decisions at \path{T030A_first_round_vote_decision} before the result is revealed at \path{T030B_first_round_result_revealed}. The same structure applies to the runoff at \path{T035A_runoff_vote_decision} and \path{T035B_final_result_revealed}. Each decision prompt includes a temporal fence that restricts the agent to information available at the simulated date.

\chapter{Tie-Handling Sensitivity Detail}
\label{app:tie_detail}

Table~\ref{tab:tie_composition} reports the composition of the top-probability ties recorded at the election decision ticks. Fifty of three hundred first-round agents and fifty-four of three hundred runoff agents produced tied top probabilities. The most common tie pattern in the first round is a three-way tie between Erdoğan, Kılıçdaroğlu, and Undecided, accounting for twenty-one agents. The most common runoff tie is a three-way tie between Erdoğan, Kılıçdaroğlu, and the Abstain/Invalid/Undecided category, accounting for twenty-five agents.

\begin{table}[!ht]
\caption{Composition of top-probability ties at election decision ticks.}
\label{tab:tie_composition}
\centering
\footnotesize
\begin{tabular}{@{}llrr@{}}
\toprule
Stage & Tied outcomes & Tie size & Agents \\
\midrule
First round & Erdoğan + Kılıçdaroğlu + Undecided                                  & 3 & 21 \\
First round & Erdoğan + Kılıçdaroğlu                                              & 2 & 13 \\
First round & Kılıçdaroğlu + Undecided                                            & 2 &  8 \\
First round & Erdoğan + Undecided                                                 & 2 &  4 \\
First round & Erdoğan + S. Oğan + Undecided                                       & 3 &  2 \\
First round & Erdoğan + S. Oğan                                                   & 2 &  1 \\
First round & Kılıçdaroğlu + Other                                                & 2 &  1 \\
\midrule
Runoff      & Abstain/Invalid/Undecided + Erdoğan + Kılıçdaroğlu                  & 3 & 25 \\
Runoff      & Abstain/Invalid/Undecided + Erdoğan                                 & 2 & 14 \\
Runoff      & Abstain/Invalid/Undecided + Kılıçdaroğlu                            & 2 &  9 \\
Runoff      & Erdoğan + Kılıçdaroğlu                                              & 2 &  6 \\
\bottomrule
\end{tabular}
\end{table}

The deterministic canonical rule resolves ties by the candidate-order specified in the aggregation code: Erdoğan, Kılıçdaroğlu, Sinan Oğan, Muharrem İnce, Other, and Undecided for the first round; Erdoğan, Kılıçdaroğlu, and Abstain/Invalid/Undecided for the runoff. The reverse-order rule inverts this ordering. The random-tie-break rule samples uniformly among the tied candidates and is averaged over ten thousand Monte Carlo iterations. Tied-agent outcomes under all three rules are reported in Table~\ref{tab:tie_sensitivity} in the main text.

\end{document}
